\documentclass[11pt]{article}

%-----------------------------------------------------------------------------%
% PACKAGES
%-----------------------------------------------------------------------------%
\usepackage[margin=1in]{geometry}
\usepackage{amsmath,amssymb,amsthm}
\usepackage{mathtools}
\usepackage{bm}
\usepackage{authblk}
\usepackage{hyperref}
\usepackage{enumitem}
\usepackage{booktabs}
\usepackage{microtype}

\hypersetup{
  colorlinks=true,
  linkcolor=blue,
  citecolor=blue,
  urlcolor=blue
}

%-----------------------------------------------------------------------------%
% THEOREM ENVIRONMENTS
%-----------------------------------------------------------------------------%
\newtheorem{definition}{Definition}[section]
\newtheorem{theorem}[definition]{Theorem}
\newtheorem{proposition}[definition]{Proposition}
\newtheorem{corollary}[definition]{Corollary}
\newtheorem{remark}[definition]{Remark}
\newtheorem{lemma}[definition]{Lemma}

%-----------------------------------------------------------------------------%
% MACROS
%-----------------------------------------------------------------------------%
\newcommand{\Z}{\mathbb{Z}}
\newcommand{\N}{\mathbb{N}}
\newcommand{\R}{\mathbb{R}}
\newcommand{\F}{\mathbb{F}}

\newcommand{\vect}[1]{\bm{#1}}
\newcommand{\norm}[1]{\left\lVert #1 \right\rVert}
\newcommand{\abs}[1]{\left\lvert #1 \right\rvert}

\DeclareMathOperator{\logit}{logit}
\DeclareMathOperator{\sigm}{sigm}

%-----------------------------------------------------------------------------%
% TITLE
%-----------------------------------------------------------------------------%
\begin{document}
\title{Multiplicity Theory Across Genomics, Cryptography, and Wearable Sensing:\\
A Prime-Indexed Framework for Certified Adaptive Systems}

\author{Citizen Gardens \\ The Prime Materia Commons}

\date{May 5, 2026}
\maketitle
\begin{abstract}
This work develops a unified, prime-indexed framework---Multiplicity Theory---for modeling nonlinear, emergent structure across genomics, cryptography, and wearable sensing as recursively stable multiplicity spaces rather than static objects. At its core, Multiplicity Theory encodes system evolution as prime-labeled interaction patterns, where multiplicity operators track the recurrence of structured events and interact with orthogonal spectral transforms and cryptographic commitments to yield scale-consistent, audit-ready dynamics. We instantiate this paradigm in three coordinated tracks.

Track~A introduces a genomic baseline for \emph{E.\ coli} K-12 that compresses 61-codon usage profiles into a six-channel Walsh--Hadamard-like (WHT) spectral representation with a GC3-controlled null ensemble. The resulting \(\Delta h\) descriptors are shown to be norm-bounded under explicit operator norms for aggregation and orthogonal mixing, enabling condition-number control and well-posed logistic regression comparisons between GC3-only and GC3+\(\Delta h\) feature sets.

Track~B develops \emph{Multiplicity-Crypto}, an integrated certification and transport layer that combines BN254 Pedersen commitments, HKDF-based directional keys, and a prime-indexed multiplicity operator over protocol transcripts. Transcript history is injected into both key derivation and associated data via multiplicity labels and hash chains, while a rank-1 coupling tensor and diagonal multiplicity operator define an encryption operator whose norm is bounded in terms of classical and quantum-like state magnitudes. A Banach-type contraction analysis yields sufficient conditions for the existence and uniqueness of fixed points in adaptive feedback regimes, providing mathematically controlled key rotation and state evolution.

Track~C extends the Codon-Aware Spectral (CAS) architecture to the WESAD wearable stress and affect detection dataset, treating multichannel physiological windows as CAS instances. High-dimensional time-series features are aggregated into a small number of semantically aligned channels and mixed via an orthogonal WHT layer, yielding CAS features with rigorously bounded operator norms. Within-subject baseline normalization and logistic regression on these CAS features inherit the same stability guarantees, and the pipeline is designed to attach directly to the Track~B multiplicity-crypto stack for model and transcript certification.

Across all three tracks, the project provides explicit linear-algebraic and operator-norm proofs, HKDF and AEAD binding constructions, and reference implementation sketches, culminating in a coherent prior art and defensive publication package. Taken together, these results demonstrate that prime-indexed multiplicity, orthogonal spectral mixing, and native cryptographic certification can be composed into a single, recursively stable design language for biological, cryptographic, and wearable-sensing systems.
\end{abstract}
\newpage
\tableofcontents
\newpage
%-----------------------------------------------------------------------------%
\section{Track A --- Genomic Baseline with 6-Channel Walsh--Hadamard Transform\\
Prior Art Defensive Publication and Technical Report}

This document specifies \emph{Track A --- Genomic Baseline}, a 2--3 week experimental protocol that applies a six-channel Walsh--Hadamard Transform (WHT) spectral decomposition to a publicly available codon-usage dataset for \emph{Escherichia coli} K-12 (approximately 4,000 coding sequences).
For each gene, a pairwise spectral distance \(\Delta h\) is computed between native codon usage and a GC3-controlled reference ensemble.
A logistic classifier is then trained on \(\Delta h\) and contrasted with a classifier trained on GC3-only features, and the absolute accuracy delta is measured.
This constitutes a Go/No-Go gate for the Track A hypothesis that WHT-based multiplicity structure captures biologically meaningful codon-usage signal beyond GC3 composition alone.
The report is written as a defensive prior-art publication to prevent third-party patent claims on the described methodology.

Track A is designed as the fastest executable testbed for Multiplicity Theory in a concrete biological domain: codon usage bias in \emph{E.~coli} K-12.\footnote{Codon usage and GC3-based analyses for \emph{E.~coli} K-12 are standard and widely used in comparative genomics; see, e.g., general codon usage tools such as GCUA and related GC3-focused pipelines, as well as public codon tables for this organism.[web:303][web:317]} 
The central question is whether a six-channel Walsh--Hadamard spectral decomposition of codon usage vectors, followed by an appropriately defined \(\Delta h\) comparison to GC3-matched null models, yields predictive structure not already explained by GC3 composition alone.

The core experimental steps are:
\begin{enumerate}[leftmargin=1.8em]
  \item Download a curated CDS set for \emph{E.~coli} K-12 (K-12 MG1655) with roughly 4,000 protein-coding genes and clean the sequences (frame-correct, multiple-of-three lengths, standard genetic code).
  \item For each gene, compute:
    \begin{itemize}
      \item a codon-usage vector \(c_g \in \R^{61}\) (excluding stops),
      \item GC content overall and at third codon positions (\(\mathrm{GC3}_g\)), and optionally other base-composition metrics.[web:303][web:304]
    \end{itemize}
  \item Apply a six-channel WHT projection mapping \(c_g\) into six interaction channels associated with multiplicity-theoretic axes (e.g., synonymous vs.\ non-synonymous, purine/pyrimidine partitions, etc.).
  \item For each gene, generate a GC3-controlled reference ensemble via codon reshuffling that preserves amino-acid sequence and GC3 profile; compute the corresponding WHT spectrum and define a per-gene \(\Delta h_g\) as the normed distance between native and reference spectra.
  \item Construct two logistic regression classifiers for a ground-truth binary label (e.g., native vs.\ foreign codon-usage status, or high vs.\ low expression class): one using only GC3-based features, and one using GC3 plus \(\Delta h_g\) (or \(\Delta h\)-derived features).
  \item Quantify the absolute accuracy delta, \(\Delta \mathrm{Acc} = \mathrm{Acc}(\text{GC3}+\Delta h) - \mathrm{Acc}(\text{GC3-only})\), using cross-validation.
\end{enumerate}

A positive \(\Delta \mathrm{Acc}\) supports the Track A thesis that \(\Delta h\) captures non-trivial, GC3-independent structure in codon usage; a null or negligible \(\Delta \mathrm{Acc}\) is a ``No-Go'' outcome that falsifies the specific spectral hypothesis at this resolution.

This document defines the mathematical framework, concrete pipeline, and representative code patterns with sufficient detail to re-implement the experiment and to establish defensive prior art.

%-----------------------------------------------------------------------------%
\subsection{Data Model and Preprocessing}
%-----------------------------------------------------------------------------%

\subsubsection{Genome and CDS Selection}

Track A assumes the use of a standard, publicly available CDS annotation for \emph{E.~coli} K-12 (e.g., MG1655) from NCBI or a similar source.
Multiple tools and datasets exist for codon usage analysis in \emph{E.~coli}, including utilities such as GCUA and coRdon that compute RSCU, GC3, and other metrics over CDS sets.[web:303][web:304][web:306]

\paragraph{Requirements.}
\begin{itemize}[leftmargin=1.8em]
  \item All CDS must be in-frame, length divisible by three, and annotated using the standard bacterial genetic code.
  \item Start and stop codons may be handled separately; in this specification stops are excluded from codon-usage vectors (yielding a 61-dimensional space).
  \item Genes shorter than a minimal length threshold (e.g., fewer than 100 codons) may be filtered out for stability.
\end{itemize}

Let \(\mathcal{G}\) denote the set of genes retained after quality control, with \(|\mathcal{G}| = G \approx 4{,}000\).

\subsubsection{Codon-Usage Vectors}

Let \(\mathcal{C}\) be the set of 61 sense codons.
For each gene \(g \in \mathcal{G}\), let \(\mathrm{count}_g(c)\) be the count of codon \(c \in \mathcal{C}\) in that gene.

\begin{definition}[Raw codon-usage vector]
The raw codon-usage vector for gene \(g\) is
\[
  u_g(c) = \mathrm{count}_g(c), \quad c \in \mathcal{C},
\]
which can be arranged as \(u_g \in \R^{61}\).
\end{definition}

\begin{definition}[Normalized codon-usage vector]
Let \(L_g = \sum_{c \in \mathcal{C}} u_g(c)\) be the codon length of gene \(g\).
The normalized codon-usage vector is
\[
  c_g(c) = \frac{u_g(c)}{L_g}, \quad c \in \mathcal{C},
\]
yielding \(c_g \in \Delta^{60}\), the 61-dimensional probability simplex.
\end{definition}

\subsubsection{GC3 and Other GC-Based Metrics}

For each gene \(g\), let \(\mathrm{GC}_g\) be the total GC fraction and \(\mathrm{GC3}_g\) the GC fraction at third codon positions.

\begin{definition}[GC3 content]
Let \(N_{3,\mathrm{G/C}}(g)\) be the number of positions at third codon sites that are G or C, and let \(N_{3,\mathrm{total}}(g)\) be the total number of third codon positions (i.e., \(L_g\)).
Then
\[
  \mathrm{GC3}_g = \frac{N_{3,\mathrm{G/C}}(g)}{N_{3,\mathrm{total}}(g)}.
\]
\end{definition}

GC3 and related metrics are standard in codon usage bias analysis, and multiple existing packages compute them directly from codon counts.[web:303][web:304][web:308]

Additional features may include overall GC, GC1, GC2, effective number of codons, or CAI, but the minimal Track A baseline focuses on GC3 (possibly plus overall GC as a secondary feature).

%-----------------------------------------------------------------------------%
\subsection{Six-Channel WHT Representation}
%-----------------------------------------------------------------------------%

\subsubsection{Walsh--Hadamard Transform Basics}

The Walsh--Hadamard Transform (WHT) provides an orthogonal decomposition over a hypercube, often used in epistasis and combinatorial genetics to decompose additive vs.\ interaction effects across binary factors.
In Track A, WHT is repurposed as a spectral transform over codon features, with six channels representing specific multiplicity-theoretic axes.

Formally, the standard WHT on \(2^n\) inputs can be written as multiplication by the \(2^n \times 2^n\) Hadamard matrix \(H^{(n)}\) with entries \(\pm 1\).
For our purposes, codons are grouped and mapped into six aggregated channels, rather than applying a full \(2^{61}\)-dimensional transform.

\subsubsection{Channel Construction}

Let there be a fixed partition of \(\mathcal{C}\) into six groups \(\mathcal{C}_1, \dots, \mathcal{C}_6\).
The grouping is chosen to reflect biologically meaningful dimensions such as:
\begin{itemize}[leftmargin=1.8em]
  \item GC3 vs.\ AT3 at synonymous sites,
  \item amino-acid hydrophobicity classes,
  \item purine vs.\ pyrimidine at first or second positions,
  \item or other structured partitions informed by multiplicity labels.
\end{itemize}
The exact grouping is part of the Track A design and must be fixed before analysis.

\begin{definition}[Channel aggregation]
Given a normalized codon-usage vector \(c_g\), define the aggregated channel intensities
\[
  a_{g,i} = \sum_{c \in \mathcal{C}_i} c_g(c), \quad i=1,\dots,6.
\]
Let \(\vect{a}_g = (a_{g,1}, \dots, a_{g,6})^\top \in \R^6\).
\end{definition}

\begin{definition}[6-channel Walsh--Hadamard transform]
Let \(H_6\) be a fixed \(6 \times 6\) orthogonal matrix with entries \(\pm 1 / \sqrt{6}\), representing a WHT-like transform on six channels (e.g., the first six rows of a normalized Hadamard matrix extended or adapted to a non-power-of-two dimension).
The 6-channel WHT spectrum for gene \(g\) is
\[
  h_g = H_6 \vect{a}_g \in \R^6.
\]
\end{definition}

\begin{remark}
The Track A protocol is agnostic to the exact choice of \(H_6\) so long as it is orthogonal, has entries of comparable magnitude, and is fixed across all genes.
The key idea is to obtain a small set of orthogonal coordinates capturing combinatorial structure across the six aggregated codon channels.
\end{remark}

%-----------------------------------------------------------------------------%
\subsection{GC3-Controlled Null Ensemble and \texorpdfstring{$\Delta h$}{Δh}}
%-----------------------------------------------------------------------------%

\subsubsection{GC3-Controlled Randomization}

To disentangle GC3-driven effects from higher-order codon structure, Track A constructs a null model of codon usage per gene that preserves:
\begin{itemize}[leftmargin=1.8em]
  \item the amino-acid sequence,
  \item the gene length,
  \item the GC3 content (up to small binning or tolerance),
\end{itemize}
while randomizing codon choices within constraints.

\begin{definition}[GC3-preserving codon reshuffling]
For gene \(g\) with amino-acid sequence \(\alpha_1,\dots,\alpha_{L_g}\) and GC3 content \(\mathrm{GC3}_g\), define a procedure that samples random synonymous codons at each position such that:
\begin{itemize}
  \item the chosen codon encodes the correct amino acid \(\alpha_i\),
  \item the overall GC3 count matches a target consistent with \(\mathrm{GC3}_g\) (exact match or small tolerance).
\end{itemize}
Repeating this sampling \(R\) times produces an ensemble of reshuffled codon-usage vectors \(\{c_{g,r}^{\mathrm{null}}\}_{r=1}^R\).
\end{definition}

Tools such as \texttt{codonuse} illustrate closely related ideas, where null codon usage is generated by randomizing codons while controlling GC3.[web:305]

\subsubsection{Null WHT Spectrum and Distance}

For each gene and each null replicate, compute the 6-channel spectrum as in the native case.

\begin{definition}[Null WHT spectra]
For gene \(g\) and null replicate \(r\), define
\[
  h_{g,r}^{\mathrm{null}} = H_6 \vect{a}_{g,r}^{\mathrm{null}} \in \R^6,
\]
where \(\vect{a}_{g,r}^{\mathrm{null}}\) is computed from \(c_{g,r}^{\mathrm{null}}\) by channel aggregation.
\end{definition}

\begin{definition}[Mean null spectrum and \(\Delta h\)]
Define the mean null spectrum
\[
  \bar{h}_g^{\mathrm{null}} = \frac{1}{R} \sum_{r=1}^R h_{g,r}^{\mathrm{null}}.
\]
Then define
\[
  \Delta h_g = \norm{h_g - \bar{h}_g^{\mathrm{null}}}_2,
\]
or more generally, the vector of component-wise deviations
\[
  \Delta \vect{h}_g = h_g - \bar{h}_g^{\mathrm{null}} \in \R^6.
\]
\end{definition}

\begin{remark}
The scalar \(\Delta h_g\) is a single summary statistic, while \(\Delta \vect{h}_g\) provides six-dimensional features.
Both can be used as inputs to downstream classifiers.
\end{remark}

%-----------------------------------------------------------------------------%
\subsection{Classification Setup and Go/No-Go Gate}
%-----------------------------------------------------------------------------%

\subsubsection{Ground-Truth Labels}

Track A requires a binary label per gene to serve as ground truth for supervised classification.
Possibilities include:
\begin{itemize}[leftmargin=1.8em]
  \item Native vs.\ foreign codon usage, as determined by an existing axis-of-variation analysis over \emph{E.~coli} K-12 genes (e.g., classification into native vs.\ horizontally transferred genes).[web:321]
  \item High vs.\ low expression classes determined from external expression data (e.g., ribosomal vs.\ lowly expressed genes).
\end{itemize}

For definiteness, let \(y_g \in \{0,1\}\) be the binary label for gene \(g\), where \(y_g = 1\) might denote ``native'' or ``highly expressed''.

\subsubsection{Feature Sets}

Define two feature sets per gene \(g\):

\paragraph{GC3-only features.}
\[
  X_g^{\mathrm{GC3}} = \bigl(\mathrm{GC3}_g\bigr) \in \R,
\]
or extended to a small GC-based vector such as \((\mathrm{GC}_g, \mathrm{GC3}_g) \in \R^2\).

\paragraph{GC3 + \(\Delta h\) features.}
\[
  X_g^{\mathrm{GC3}+\Delta h} =
  \begin{cases}
    \bigl(\mathrm{GC3}_g, \Delta h_g\bigr) \in \R^2, & \text{scalar \(\Delta h\) case},\\[0.5em]
    \bigl(\mathrm{GC3}_g, \Delta h_{g,1}, \dots, \Delta h_{g,6}\bigr) \in \R^7, & \text{vector \(\Delta \vect{h}_g\) case}.
  \end{cases}
\]

\subsubsection{Logistic Regression Models}

\begin{definition}[GC3-only logistic model]
Let \(\theta^{\mathrm{GC3}} \in \R^{d_1}\) with \(d_1 = 1\) or \(2\) depending on the GC3 feature dimension.
The GC3-only logistic model is
\[
  \Pr(y_g = 1 \mid X_g^{\mathrm{GC3}}) =
  \sigma\bigl((\theta^{\mathrm{GC3}})^\top X_g^{\mathrm{GC3}}\bigr),
  \quad \sigma(z) = \frac{1}{1+e^{-z}}.
\]
\end{definition}

\begin{definition}[GC3 + \(\Delta h\) logistic model]
Let \(\theta^{\mathrm{GC3}+\Delta h} \in \R^{d_2}\) with \(d_2 = 2\) or \(7\).
The expanded logistic model is
\[
  \Pr(y_g = 1 \mid X_g^{\mathrm{GC3}+\Delta h}) =
  \sigma\bigl((\theta^{\mathrm{GC3}+\Delta h})^\top X_g^{\mathrm{GC3}+\Delta h}\bigr).
\]
\end{definition}

These models can be trained via standard \(\ell_2\)-regularized maximum likelihood or stochastic gradient methods, as commonly used for sequence-based logistic classifiers.[web:307]

\subsubsection{Accuracy and Absolute Delta}

For each model, evaluate accuracy via cross-validation:

\begin{definition}[Cross-validated accuracy]
Partition \(\mathcal{G}\) into \(K\) folds.
For each fold, train the model on \(K-1\) folds and evaluate on the held-out fold, then average.
Denote the resulting accuracies as \(\mathrm{Acc}_{\mathrm{GC3}}\) and \(\mathrm{Acc}_{\mathrm{GC3}+\Delta h}\).
\end{definition}

\begin{definition}[Absolute accuracy delta]
The Track A Go/No-Go metric is
\[
  \Delta \mathrm{Acc} = \mathrm{Acc}_{\mathrm{GC3}+\Delta h} - \mathrm{Acc}_{\mathrm{GC3}}.
\]
\end{definition}

\begin{remark}[Go/No-Go interpretation]
A substantially positive \(\Delta \mathrm{Acc}\) (beyond noise and confidence intervals) indicates that the WHT-derived \(\Delta h\) feature carries predictive information not captured by GC3 alone, supporting the Track A spectral hypothesis.
A negligible or zero \(\Delta \mathrm{Acc}\) is a ``No-Go'' outcome, suggesting that codon-usage signal relevant to the chosen labels is well-approximated by GC3 composition and that the proposed 6-channel WHT does not add measurable discriminative power under this design.
\end{remark}

%-----------------------------------------------------------------------------%
\subsection{Representative Code Snippet (Python)}
%-----------------------------------------------------------------------------%

This section provides representative pseudocode in Python (using NumPy and scikit-learn) for the Track A pipeline.
It is not meant to be a drop-in implementation but rather a concrete sketch of the key steps.

\subsubsection{Feature Extraction and \texorpdfstring{$\Delta h$}{Δh}}

\begin{verbatim}
import numpy as np
from sklearn.linear_model import LogisticRegression
from sklearn.model_selection import StratifiedKFold
from Bio import SeqIO  # from Biopython

CODONS = [  # 61 sense codons, stops excluded
    "AAA","AAC","AAG","AAT", # ...
    # (fill with all codons except TAA, TAG, TGA)
]

CODON_IDX = {c: i for i, c in enumerate(CODONS)}

# Example 6x6 orthogonal transform (placeholder – must be fixed beforehand)
H6 = (1.0 / np.sqrt(6)) * np.array([
    [ 1,  1,  1,  1,  1,  1],
    [ 1, -1,  1, -1,  1, -1],
    [ 1,  1, -1, -1,  1,  1],
    [ 1, -1, -1,  1,  1, -1],
    [ 1,  1,  1, -1, -1, -1],
    [ 1, -1,  1,  1, -1,  1],
])

# User-defined mapping of codons to 6 channels
def codon_channel(codon: str) -> int:
    # TODO: define biologically meaningful partition
    # For now, a placeholder mapping based on GC3 etc.
    # Return an integer 0..5.
    raise NotImplementedError

def gc3_content(seq: str) -> float:
    thirds = seq[2::3]
    gc = sum(1 for b in thirds if b in "GCgc")
    return gc / max(len(thirds), 1)

def codon_usage_vector(seq: str) -> np.ndarray:
    counts = np.zeros(len(CODONS), dtype=float)
    for i in range(0, len(seq) - 2, 3):
        codon = str(seq[i:i+3]).upper()
        if codon in CODON_IDX:
            counts[CODON_IDX[codon]] += 1
    total = counts.sum()
    if total > 0:
        counts /= total
    return counts

def six_channel_projection(codon_vec: np.ndarray) -> np.ndarray:
    channel_sums = np.zeros(6, dtype=float)
    for codon, idx in CODON_IDX.items():
        ch = codon_channel(codon)
        channel_sums[ch] += codon_vec[idx]
    return H6 @ channel_sums  # 6D spectrum

def random_gc3_preserving_sequence(seq: str, target_gc3: float) -> str:
    # Placeholder; real implementation must enforce amino acid and GC3 constraints.
    raise NotImplementedError

def compute_delta_h_for_gene(seq: str, R: int = 100) -> float:
    # Native spectrum
    cu_native = codon_usage_vector(seq)
    h_native = six_channel_projection(cu_native)

    # GC3-controlled null ensemble
    gc3 = gc3_content(seq)
    h_null = []
    for _ in range(R):
        seq_null = random_gc3_preserving_sequence(seq, gc3)
        cu_null = codon_usage_vector(seq_null)
        h_null.append(six_channel_projection(cu_null))
    h_null = np.vstack(h_null)
    h_null_mean = h_null.mean(axis=0)

    delta = np.linalg.norm(h_native - h_null_mean)
    return float(delta)
\end{verbatim}

\subsubsection{Logistic Regression and Accuracy Delta}

\begin{verbatim}
def run_trackA_classification(cds_fasta_path, labels_dict, n_splits=5):
    seqs = []
    gc3_vals = []
    delta_h_vals = []
    ys = []

    for rec in SeqIO.parse(cds_fasta_path, "fasta"):
        gene_id = rec.id
        if gene_id not in labels_dict:
            continue
        seq = str(rec.seq)
        if len(seq) % 3 != 0:
            continue  # enforce in-frame

        gc3 = gc3_content(seq)
        d_h = compute_delta_h_for_gene(seq)

        seqs.append(seq)
        gc3_vals.append(gc3)
        delta_h_vals.append(d_h)
        ys.append(labels_dict[gene_id])

    gc3_vals = np.array(gc3_vals).reshape(-1, 1)
    delta_h_vals = np.array(delta_h_vals).reshape(-1, 1)
    ys = np.array(ys)

    X_gc3 = gc3_vals
    X_gc3_dh = np.hstack([gc3_vals, delta_h_vals])

    clf_gc3 = LogisticRegression(max_iter=1000)
    clf_gc3_dh = LogisticRegression(max_iter=1000)

    skf = StratifiedKFold(n_splits=n_splits, shuffle=True, random_state=42)

    acc_gc3, acc_gc3_dh = [], []

    for train_idx, test_idx in skf.split(X_gc3, ys):
        X_tr_gc3, X_te_gc3 = X_gc3[train_idx], X_gc3[test_idx]
        X_tr_dh,  X_te_dh  = X_gc3_dh[train_idx], X_gc3_dh[test_idx]
        y_tr, y_te = ys[train_idx], ys[test_idx]

        clf_gc3.fit(X_tr_gc3, y_tr)
        clf_gc3_dh.fit(X_tr_dh, y_tr)

        acc_gc3.append(clf_gc3.score(X_te_gc3, y_te))
        acc_gc3_dh.append(clf_gc3_dh.score(X_te_dh, y_te))

    acc_gc3 = float(np.mean(acc_gc3))
    acc_gc3_dh = float(np.mean(acc_gc3_dh))
    delta_acc = acc_gc3_dh - acc_gc3

    return acc_gc3, acc_gc3_dh, delta_acc
\end{verbatim}

This code illustrates the minimal Track A computation of \(\Delta h\) per gene and the comparison of GC3-only vs.\ GC3+\(\Delta h\) logistic models under cross-validation.

%-----------------------------------------------------------------------------%
\subsection{Implementation Timeline and Practical Notes}
%-----------------------------------------------------------------------------%

\subsubsection{Timeline (2--3 Weeks)}

A realistic 2--3 week schedule for Track A is:

\begin{itemize}[leftmargin=1.8em]
  \item \textbf{Week 1:} Acquire and clean CDS for \emph{E.~coli} K-12, implement codon counting and GC3 computation, implement 6-channel mapping and WHT projection.
  \item \textbf{Week 2:} Implement GC3-preserving randomization, compute \(\Delta h\) per gene, assemble labels, and run initial logistic regression experiments.
  \item \textbf{Week 3 (optional):} Refine channel design, explore alternative labels (e.g., native vs.\ foreign axes), and perform robustness checks (different randomization schemes, different WHT bases, etc.).
\end{itemize}

\subsubsection{Design Choices and Extensions}

Key design choices that can be varied in follow-up experiments include:

\begin{itemize}[leftmargin=1.8em]
  \item Alternative channel partitions \(\mathcal{C}_i\) (e.g., focusing more tightly on GC3-synonymous groups).
  \item Using full \(\Delta \vect{h}_g\) (6 channels) rather than scalar \(\Delta h_g\).
  \item Incorporating additional GC-based features, RSCU, or CAI as baselines.
  \item Switching to other linear classifiers (e.g., SVMs or elastic-net logistic regression) while preserving the GC3 vs.\ GC3+\(\Delta h\) comparison.
\end{itemize}

%-----------------------------------------------------------------------------%
\subsection{Defensive Publication Statement}
%-----------------------------------------------------------------------------%

\subsubsection*{Scope}

This document constitutes a defensive publication covering the following Track A contributions:

\begin{enumerate}[leftmargin=1.8em]
  \item The use of a six-channel WHT spectral decomposition over codon-usage vectors for genome-wide analysis of \emph{E.~coli} K-12 genes.
  \item The definition of a GC3-preserving null ensemble per gene and the resulting per-gene \(\Delta h\) as the normed difference between native and null WHT spectra.
  \item The use of \(\Delta h\) as a covariate in a logistic regression classifier alongside GC3 (and related GC-based features) to test for non-GC3 codon-usage signal.
  \item The Go/No-Go gate defined in terms of the absolute accuracy delta between GC3-only and GC3+\(\Delta h\) models in supervised classification of gene labels.
\end{enumerate}

\subsubsection*{Intent}

By publishing this specification with technical and mathematical detail sufficient for independent re-implementation, the authors intend to create prior art that prevents third parties from obtaining patent rights over implementations of Track A as described here, including variants that differ only in straightforward parameter tuning or channel selection.


\section{Track B: Multiplicity-Crypto Integrated Certification and Transport}


This document specifies \emph{Track B}, a multiplicity-theoretic certification and transport framework built on top of the native Multiplicity Crypto stack.
Track B binds stateful data products to BN254 Pedersen commitments, prime-indexed multiplicity operators, and transcript-bound HKDF-derived AEAD keys, with optional Quantum-State-Based Security (QSBS) extensions and QKD hybrid key substrates.
The goal of this report is twofold:
(1) to provide a mathematically complete specification suitable for independent re-implementation, and
(2) to constitute prior art for defensive patent purposes, preventing third-party claims on the techniques described herein.

%-----------------------------------------------------------------------------%
\subsection{Executive Summary}
%-----------------------------------------------------------------------------%

\subsubsection{Purpose of Track B}

Track B is defined as a second, cryptography-focused track in the PhaseMirror / Multiplicity ecosystem whose primary role is to \emph{certify} and \emph{securely transport} evolving data products and model states.
Rather than treating Track B as a simple logging or checksum layer, the design interprets Track B as a multiplicity-governed protocol:
each transformation, message, and acknowledgement becomes a prime-labeled interaction step whose history is encoded and preserved in a multiplicity operator.

Concretely, Track B integrates:
\begin{itemize}[leftmargin=1.5em]
  \item A BN254-based Pedersen commitment engine (via Multiplicity Crypto) for circuit-native commitments.
  \item A prime-indexed multiplicity operator that G\"odel-encodes Track B transcripts into uniquely factorizable integers.
  \item A transcript hash chain and context hash used as associated data for authenticated encryption.
  \item An HKDF-based key schedule with multiplicity labels embedded into the info string, producing directional AEAD keys.
  \item An optional QSBS overlay that maps classical payloads and quantum-style amplitudes into a joint frequency space and couples them through a low-rank tensor.
\end{itemize}

\subsubsection{Defensive Publication Scope}

This document is intended as a defensive prior-art publication and covers at least the following contributions related to Track B:

\begin{enumerate}[leftmargin=1.5em,label=\textbf{B.\arabic*}]
  \item \textbf{Track B state certification} via BN254 Pedersen commitments over frequency-mapped state representations and multiplicity-labeled transcripts.
  \item \textbf{Multiplicity-indexed HKDF derivation for Track B} directional keys, including an 8-byte multiplicity label embedded in the info string.
  \item \textbf{Frequency-unified state commitments for Track B}, where classical and (optionally) quantum-like state descriptors are combined into a single commitment handle.
  \item \textbf{Prime-indexed multiplicity transcript binding for Track B}, giving recursively stable transcript identities across Track B turns.
  \item \textbf{Adaptive Track B key rotation governed by feedback dynamics} on error, latency, and load metrics, expressed in a multiplicity-theoretic state space.
\end{enumerate}

All such techniques are hereby described with sufficient detail for reconstruction and are intended to be irrevocably placed into the public domain for defensive purposes.

%-----------------------------------------------------------------------------%
\subsection{Track B Conceptual Architecture}
%-----------------------------------------------------------------------------%

\subsubsection{Layered View}

Track B is organized as a layered system, specialized for certification and transport:

\begin{itemize}[leftmargin=1.5em]
  \item \textbf{B0 -- Curve and primitives:} BN254 group and scalar field, random oracle, AEAD scheme.
  \item \textbf{B1 -- Commitment layer:} BN254 Pedersen commitments for Track B states.
  \item \textbf{B2 -- Frequency mapping:} maps Track B states into joint frequency vectors.
  \item \textbf{B3 -- Multiplicity and transcripts:} prime-indexed transcript encoding and multiplicity operator.
  \item \textbf{B4 -- Key derivation:} HKDF-based directional keys with multiplicity labels.
  \item \textbf{B5 -- Encryption/AEAD:} AEAD encryption with transcript-bound associated data.
  \item \textbf{B6 -- Feedback dynamics:} adaptive key rotation and state evolution controlled by observed metrics.
  \item \textbf{B7 -- Optional QSBS overlay:} joint classical/quantum frequency model integrated into the above.
\end{itemize}

Each layer is designed to be independently testable but composable, with clear functional boundaries and simple interfaces.

\subsubsection{Track B State and Roles}

Track B is parameterized by a discrete time index \(t \in \N\) and involves at least two roles (e.g., producer and consumer, or module A and module B).
Its instantaneous state is:

\[
  S_t^{(B)} = \bigl(x_t^{(B)}, \sigma_t^{(B)}, M_t^{(B)}, T_t^{(B)}\bigr),
\]

optionally extended with a quantum or amplitude-like component \(|\psi_t^{(B)}\rangle\) when QSBS is used.

\begin{itemize}[leftmargin=1.5em]
  \item \(x_t^{(B)}\): classical Track B payload (bitstring or structured data serialized to bits).
  \item \(\sigma_t^{(B)}\): Track B transcript (ordered message history and labels).
  \item \(M_t^{(B)}\): multiplicity operator, a vector of prime-indexed message-class counts.
  \item \(T_t^{(B)}\): coupling tensor, modeling interactions between classical and quantum modes.
\end{itemize}

The goal is that every emission and transformation of Track B is both cryptographically committed and multiplicity-identifiable.

%-----------------------------------------------------------------------------%
\subsection{Mathematical Overview}
%-----------------------------------------------------------------------------%

\subsubsection{Standing Assumptions}

\paragraph{Cryptographic primitives.}
We assume:
\begin{itemize}[leftmargin=1.5em]
  \item A pairing-friendly elliptic curve BN254 with group \(\G_1\) of prime order \(r\) and generator \(G\).
  \item A random-oracle hash function \(\SHA: \{0,1\}^* \to \{0,1\}^{256}\).
  \item An HMAC-based HKDF key derivation function.
  \item A nonce-based AEAD scheme (e.g., AES-256-GCM) that is IND-CPA and INT-CTXT secure.
\end{itemize}

\paragraph{Security parameter.}
Let \(\lambda\) denote the security parameter; all polynomial-time bounds are with respect to \(\lambda\).
Negligible functions are denoted \(\mathsf{negl}(\lambda)\).

\subsubsection{Track B Pedersen Commitments}

\begin{definition}[BN254 Pedersen commitment for Track B]
Let \(\G_1\) be the prime-order subgroup of BN254 with generator \(G\), and \(\F_r\) the scalar field of order \(r\).
Define an independent generator \(H \in \G_1\) as
\[
  H = h_G \cdot G, \quad
  h_G = \bigl(\mathrm{int}(\SHA(\texttt{"pedersen\_h"} \| \mathrm{enc}(G))) \bmod r\bigr) \in \F_r.
\]
For a Track B value \(v \in \F_r\) and blinding factor \(\rho \in \F_r\), the commitment is
\[
  C^{(B)}(v, \rho) = \rho G + v H \in \G_1.
\]
\end{definition}

This is a standard Pedersen commitment specialized to BN254 and used as the canonical commitment of Track B state encodings.

\subsubsection{Prime-Indexed Transcript Encoding and Multiplicity}

Let \(\{p_i\}_{i \ge 1}\) denote the increasing sequence of primes.

\begin{definition}[Track B message classes]
Let \(\mathcal{C} = \{c_1, \dots, c_k\}\) be the set of distinct Track B message classes (e.g., \texttt{state\_update}, \texttt{ack}, \texttt{error}, etc.).
For a transcript \(\sigma_t^{(B)}\), let \(c_i(\sigma_t^{(B)}) \in \N\) be the count of messages of class \(c_i\) up to time \(t\).
\end{definition}

\begin{definition}[Prime-indexed transcript encoding]
The Track B transcript encoding is the integer
\[
  \mathrm{enc}_B(\sigma_t^{(B)}) = \prod_{i=1}^k p_i^{c_i(\sigma_t^{(B)})}.
\]
\end{definition}

\begin{definition}[Track B multiplicity operator]
The Track B multiplicity operator at time \(t\) is the vector
\[
  M_t^{(B)} = \bigl(c_1(\sigma_t^{(B)}), \dots, c_k(\sigma_t^{(B)})\bigr) \in \Z_{\ge 0}^k.
\]
\end{definition}

\begin{proposition}[Injectivity of Track B transcript encoding]
Distinct Track B message-class multisets yield distinct encodings.
Equivalently, the map
\(
\sigma_t^{(B)} \mapsto \mathrm{enc}_B(\sigma_t^{(B)})
\)
is injective on multisets of classes.
\end{proposition}

\begin{proof}
By the Fundamental Theorem of Arithmetic, the factorization of a positive integer into primes is unique.
Since each message class \(c_i\) is assigned its own prime \(p_i\), the exponents of \(p_i\) in \(\mathrm{enc}_B(\sigma_t^{(B)})\) are exactly \(c_i(\sigma_t^{(B)})\).
Thus the multiset \(\bigl\{(c_i, c_i(\sigma_t^{(B)}))\bigr\}\) is uniquely determined by the encoding.
\end{proof}

\begin{theorem}[Operator norms of \(M_t^{(B)}\)]
Let \(N_t = \sum_{i=1}^k c_i(\sigma_t^{(B)})\).
Then:
\begin{align*}
  \norm{M_t^{(B)}}_1 &= N_t, \\
  \norm{M_t^{(B)}}_2 &\le N_t, \\
  \norm{M_t^{(B)}}_\infty &\le N_t.
\end{align*}
\end{theorem}

\begin{proof}
The \(\ell_1\) norm of \(M_t^{(B)}\) is \(\sum_i c_i = N_t\).
Since each component is non-negative and bounded by \(N_t\), we have
\(
\norm{M_t^{(B)}}_2^2 = \sum_i c_i^2 \le \left(\max_i c_i\right)\sum_i c_i \le N_t^2
\),
so \(\norm{M_t^{(B)}}_2 \le N_t\).
Similarly, \(\norm{M_t^{(B)}}_\infty = \max_i c_i \le N_t\).
\end{proof}

\subsubsection{Frequency Mapping for Track B}

Track B uses a frequency map to embed its state into a Euclidean space suitable for unified commitments.

\begin{definition}[Classical Track B frequency map]
Fix two distinct frequencies \(f_0, f_1 \in \R\).
Let \(x_t^{(B)} \in \{0,1\}^n\) be the bitstring encoding of the Track B payload at time \(t\).
Define the classical frequency map
\[
  F_c^{(B)}(x_t^{(B)}) = (f_{c,1}, \dots, f_{c,n}) \in \R^n,
\]
where
\[
  f_{c,i} =
  \begin{cases}
    f_0 & \text{if } (x_t^{(B)})_i = 0, \\
    f_1 & \text{if } (x_t^{(B)})_i = 1.
  \end{cases}
\]
\end{definition}

Optionally, if Track B incorporates quantum-like or probabilistic amplitude vectors, we can define:

\begin{definition}[Quantum-like Track B frequency map]
Let \(|\psi_t^{(B)}\rangle = \sum_{j=1}^m \alpha_j |j\rangle\) with complex amplitudes \(\alpha_j\).
Define
\[
  F_q^{(B)}(|\psi_t^{(B)}\rangle) = 
  \bigl(|\alpha_1|, \arg \alpha_1, \dots, |\alpha_m|, \arg \alpha_m\bigr) \in \R^{2m},
\]
where \(\arg\) is taken in a fixed principal range (e.g., \([0,2\pi)\)).
\end{definition}

\begin{definition}[Joint Track B frequency vector]
The joint frequency representation of Track B state at time \(t\) is
\[
  F_t^{(B)} =
  \begin{cases}
    F_c^{(B)}(x_t^{(B)}) \in \R^n & \text{(purely classical case)},\\[0.5em]
    \bigl(F_c^{(B)}(x_t^{(B)}), F_q^{(B)}(|\psi_t^{(B)}\rangle)\bigr) \in \R^{n+2m} & \text{(QSBS case)}.
  \end{cases}
\]
\end{definition}

\subsubsection{Coupling Tensor and Encryption Operator}

In the QSBS-enabled version of Track B, a coupling tensor captures interactions between classical and amplitude-mode components.

\begin{definition}[Track B coupling tensor]
Let \(u_t = F_c^{(B)}(x_t^{(B)}) \in \R^n\) and \(v_t = F_q^{(B)}(|\psi_t^{(B)}\rangle) \in \R^{2m}\).
Define the rank-1 coupling tensor
\[
  T_t^{(B)} = u_t v_t^\top \in \R^{n \times 2m}.
\]
\end{definition}

\begin{proposition}[Norms of \(T_t^{(B)}\)]
Let \(B_c = \sup_t \norm{u_t}_2\) and \(B_q = \sup_t \norm{v_t}_2\).
Then:
\begin{align*}
  \norm{T_t^{(B)}}_F &= \norm{u_t}_2 \norm{v_t}_2 \le B_c B_q, \\
  \norm{T_t^{(B)}}_{2\to 2} &= \norm{u_t}_2 \norm{v_t}_2 \le B_c B_q.
\end{align*}
\end{proposition}

\begin{proof}
For a rank-1 matrix \(T = u v^\top\), the Frobenius norm is
\(
\norm{T}_F^2 = \sum_{ij} u_i^2 v_j^2 = \norm{u}_2^2 \norm{v}_2^2
\),
and the only non-zero singular value is \(\sigma_1 = \norm{u}_2\norm{v}_2\).
Thus \(\norm{T}_F = \norm{T}_{2\to 2} = \norm{u}_2\norm{v}_2\).
The bound follows by definition of \(B_c,B_q\).
\end{proof}

\begin{definition}[Track B encryption operator]
Let \(S_t^{(B)} = (x_t^{(B)}, \sigma_t^{(B)}, M_t^{(B)}, T_t^{(B)})\) and let \(F_B\) be a (possibly non-linear) feedback function.
The Track B encryption operator is
\[
  \EncOp_B(t) = M_t^{(B)} \bigl( T_t^{(B)} S_t^{(B)} \bigr) + F_B(S_t^{(B)}).
\]
\end{definition}

The exact action of \(M_t^{(B)}\) on \(T_t^{(B)} S_t^{(B)}\) depends on how the state is embedded into a vector space;
for the purpose of analysis, it suffices to note that the resulting vector norm is bounded in terms of \(N_t, B_c, B_q\) and bounds on \(F_B\).

\begin{proposition}[Norm bound on Track B encryption output]
Assume there exists \(D_S, D_F > 0\) such that \(\norm{S_t^{(B)}}_2 \le D_S\) and \(\norm{F_B(S_t^{(B)})}_2 \le D_F\) for all \(t\).
Then
\[
  \norm{\EncOp_B(t)}_2 \le N_t B_c B_q D_S + D_F.
\]
\end{proposition}

\begin{proof}
Using the standard inequalities:
\[
  \norm{\EncOp_B(t)}_2 \le \norm{M_t^{(B)} T_t^{(B)} S_t^{(B)}}_2 + \norm{F_B(S_t^{(B)})}_2.
\]
Treating \(M_t^{(B)}\) as a diagonal operator with spectral norm bounded by \(N_t\), and using the operator norms for \(T_t^{(B)}\),
\[
  \norm{M_t^{(B)} T_t^{(B)} S_t^{(B)}}_2 \le \norm{M_t^{(B)}}_{2\to 2} \norm{T_t^{(B)}}_{2\to 2} \norm{S_t^{(B)}}_2 \le N_t B_c B_q D_S.
\]
Adding the bound on \(F_B\) yields the stated inequality.
\end{proof}

\subsubsection{Feedback Dynamics and Contraction}

\begin{definition}[Track B feedback map]
Let \(\Theta_t = (\varepsilon_t, \lambda_t, \mathrm{QBER}_t, \ell_t)\) denote the vector of observed metrics at time \(t\) (e.g., error rate, latency, quantum bit error rate, system load).
The Track B feedback map is
\[
  f_B : (M_t^{(B)}, T_t^{(B)}, \Theta_t) \mapsto (M_{t+1}^{(B)}, T_{t+1}^{(B)}).
\]
\end{definition}

\begin{definition}[Product metric space]
Equip \(\R^k \times \R^{n\times 2m}\) with the norm
\[
  \norm{(M,T)} := \norm{M}_2 + \norm{T}_F.
\]
\end{definition}

\begin{theorem}[Banach contraction for Track B feedback]
Suppose there exist non-negative constants \(L_M, L_T\) such that for all \(M,T,M',T'\),
\begin{align*}
  \norm{f_B(M,T,\Theta)_M - f_B(M',T',\Theta)_M}_2 &\le L_M \norm{(M,T)-(M',T')}, \\
  \norm{f_B(M,T,\Theta)_T - f_B(M',T',\Theta)_T}_F &\le L_T \norm{(M,T)-(M',T')},
\end{align*}
and let \(L = L_M + L_T\).
If \(L < 1\) for the regime of interest, then:
\begin{enumerate}[leftmargin=1.5em]
  \item There exists a unique fixed point \((M_\ast^{(B)}, T_\ast^{(B)})\).
  \item For any initialization \((M_0^{(B)}, T_0^{(B)})\), the iterates converge geometrically:
  \[
    \norm{(M_t^{(B)}, T_t^{(B)}) - (M_\ast^{(B)}, T_\ast^{(B)})} \le L^t \norm{(M_0^{(B)}, T_0^{(B)}) - (M_\ast^{(B)}, T_\ast^{(B)})}.
  \]
\end{enumerate}
\end{theorem}

\begin{proof}
The space \(\R^k \times \R^{n\times 2m}\) equipped with the given norm is complete.
The Lipschitz assumptions imply
\[
  \norm{f_B(M,T,\Theta) - f_B(M',T',\Theta)} \le (L_M + L_T) \norm{(M,T)-(M',T')} = L \norm{(M,T)-(M',T')}.
\]
Thus \(f_B\) is a contraction with constant \(L < 1\), and the Banach Fixed-Point Theorem applies.
\end{proof}

This provides a mathematical basis for eventual stabilisation of multiplicity labels used in Track B HKDF derivation and for designing key-rotation policies.

%-----------------------------------------------------------------------------%
\subsection{Track B Key Derivation and Encryption}
%-----------------------------------------------------------------------------%

\subsubsection{Transcript Chains and Context Hash}

Track B uses per-role transcript chains and a combined context hash for binding keys and associated data to the protocol history.

\begin{definition}[Per-role transcript chain]
For each role \(s \in \{A,B\}\), define the chain recursively:
\[
  \mathrm{chain}_{s,0} = 0^{256}, \qquad
  \mathrm{chain}_{s,t} = \SHA\bigl(\mathrm{chain}_{s,t-1} \, \| \, m_{s,t}\bigr),
\]
where \(m_{s,t}\) is the \(t\)-th message from sender \(s\) in Track B.
\end{definition}

\begin{definition}[Track B context hash]
The Track B context hash at time \(t\) is
\[
  \mathrm{ctx}_t^{(B)} = \SHA\bigl(\mathrm{chain}_{A,t} \, \| \, \mathrm{chain}_{B,t}\bigr).
\]
\end{definition}

\subsubsection{Multiplicity-Indexed HKDF Derivation}

\begin{definition}[Multiplicity label for Track B]
Let \(M_t^{(B)}\) be the multiplicity vector for Track B at time \(t\).
Define an 8-byte encoding
\[
  \mathrm{mult}_t^{(B)} = \mathrm{Enc}_8\bigl(M_t^{(B)}\bigr),
\]
where \(\mathrm{Enc}_8\) is any injective encoding from the relevant range into 8 bytes (e.g., by hashing or truncating a canonical encoding).
\end{definition}

\begin{definition}[Directional Track B AEAD key]
Let \(K_t\) be a shared secret (e.g., from QKD or an existing key-agreement scheme),
let \(\mathrm{salt}_t\) be a salt, and let \(d \in \{A\to B, B\to A\}\) denote direction.
Define the info string
\[
  \mathrm{info}_d^{(B)} = \texttt{"QSBS-v1:enc:"} \, \| \, d \, \| \, \mathrm{ctx}_t^{(B)} \, \| \, \mathrm{mult}_t^{(B)},
\]
and the directional AEAD key
\[
  K_{\mathrm{enc},d,t}^{(B)} = \HKDF\bigl(K_t; \mathrm{salt}_t, \mathrm{info}_d^{(B)}\bigr).
\]
\end{definition}

This design treats the multiplicity label as authenticated meta-information influencing key derivation, rather than as a source of entropy, preserving compatibility with standard HKDF security analyses.

\subsubsection{Track B AEAD Encryption}

\begin{definition}[Track B ciphertext]
Let \(x_t^{(B)}\) be the plaintext payload, let \(\mathrm{nonce}_{d,t}\) be a nonce,
and let \(\mathrm{ad}_t^{(B)}\) be associated data defined as
\[
  \mathrm{ad}_t^{(B)} = \bigl(\mathrm{ctx}_t^{(B)}, F_t^{(B)}, M_t^{(B)}\bigr)
\]
encoded in some deterministic, unambiguous format.
The Track B ciphertext is
\[
  C_t^{(B)} = \Enc_{K_{\mathrm{enc},d,t}^{(B)}}\bigl(\mathrm{nonce}_{d,t}, x_t^{(B)}, \mathrm{ad}_t^{(B)}\bigr).
\]
\end{definition}

Decryption proceeds by recomputing the same key and associated data from the transcript and multiplicity state, then verifying the AEAD tag.

%-----------------------------------------------------------------------------%
\subsection{Track B State Certification}
%-----------------------------------------------------------------------------%

\subsubsection{Canonical Commitment Handle}

\begin{definition}[Track B commitment handle]
Given a Track B state \(S_t^{(B)}\), compute the frequency representation \(F_t^{(B)}\)
and serialize it to a string \(f_t\) via a deterministic encoding (e.g., fixed-precision decimal).
Let \(v_t \in \F_r\) be the reduction of \(\SHA(f_t)\) modulo \(r\) and let \(\rho_t\) be a scalar derived via a domain-separated hash.
Define
\[
  h_t^{(B)} = C^{(B)}(v_t, \rho_t).
\]
\end{definition}

This handle is used as the canonical Track B commitment to the state at time \(t\).
It can be published, stored, or verified in a ZK circuit without revealing the underlying state.

\subsubsection{Audit and Verification}

Given a claimed Track B state \(S_t^{(B)}\) and blinding \(\rho_t\), verification consists of:
\begin{enumerate}[leftmargin=1.5em]
  \item Recomputing \(F_t^{(B)}\) and its serialized representation \(f_t\).
  \item Computing \(v_t = \mathrm{int}(\SHA(f_t)) \bmod r\).
  \item Computing \(C^{(B)}(v_t,\rho_t)\) and checking equality to the published handle \(h_t^{(B)}\).
\end{enumerate}

This can be implemented either off-chain (e.g., in Rust/TypeScript) or inside a ZK circuit over BN254.

%-----------------------------------------------------------------------------%
\subsection{Example Code Snippet}
%-----------------------------------------------------------------------------%

This section provides representative (non-executable in this document) code snippets illustrating a possible implementation of the core Track B primitives using a BN254/WASM stack and a TypeScript wrapper.

\subsubsection{BN254 Pedersen Commitment (Rust)}

\begin{verbatim}
use ark_bn254::{G1Projective, Fr};
use ark_ec::Group;
use ark_ff::{PrimeField, BigInteger};
use sha2::{Sha256, Digest};
use wasm_bindgen::prelude::*;

/// derive_scalar(d, s) = SHA-256(d || s) mod r
fn derive_scalar(domain: &str, data: &str) -> Fr {
    let mut h = Sha256::new();
    h.update(domain.as_bytes());
    h.update(data.as_bytes());
    let bytes = h.finalize();
    Fr::from_be_bytes_mod_order(&bytes)
}

/// H = SHA-256("pedersen_h" || enc(G)) * G
fn pedersen_generator_h() -> G1Projective {
    let g = G1Projective::generator();
    let h_scalar = derive_scalar("pedersen_h", "generator");
    g * h_scalar
}

/// Compute Track B commitment handle over BN254.
/// Returns x-coordinate in hex for circuit interop.
#[wasm_bindgen]
pub fn trackb_commitment(serialized_state: &str, salt: &str) -> String {
    let g = G1Projective::generator();
    let h = pedersen_generator_h();
    // Reduce SHA-256(serialized_state) mod r for v
    let v = derive_scalar("trackb_value", serialized_state);
    // Domain-separated blinding
    let rho = derive_scalar("trackb_blinding", salt);
    let commitment = (g * rho + h * v).into_affine();
    let x_bytes = commitment.x.into_bigint().to_bytes_be();
    format!("0x{}", hex::encode(&x_bytes))
}
\end{verbatim}

\subsubsection{Transcript and Key Derivation (TypeScript)}

\begin{verbatim}
import { createHash, createHmac } from "crypto";
import init, { trackb_commitment } from "./wasm/multiplicity_crypto_wasm";

type Direction = "A2B" | "B2A";

export class TrackBTranscript {
  private chainA: Buffer;
  private chainB: Buffer;

  constructor() {
    this.chainA = Buffer.alloc(32, 0);
    this.chainB = Buffer.alloc(32, 0);
  }

  update(role: "A" | "B", message: Buffer) {
    const prev = role === "A" ? this.chainA : this.chainB;
    const next = createHash("sha256").update(prev).update(message).digest();
    if (role === "A") this.chainA = next;
    else this.chainB = next;
  }

  contextHashHex(): string {
    const ctx = createHash("sha256")
      .update(this.chainA)
      .update(this.chainB)
      .digest("hex");
    return "0x" + ctx;
  }

  chains(): { chainA: Buffer; chainB: Buffer } {
    return { chainA: Buffer.from(this.chainA), chainB: Buffer.from(this.chainB) };
  }
}

function hkdfExtract(salt: Buffer, ikm: Buffer): Buffer {
  return createHmac("sha256", salt).update(ikm).digest();
}

function hkdfExpand(prk: Buffer, info: Buffer, length: number): Buffer {
  const okm = Buffer.alloc(length);
  let t = Buffer.alloc(0);
  let offset = 0;
  for (let i = 1; offset < length; i++) {
    t = createHmac("sha256", prk)
      .update(t)
      .update(info)
      .update(Buffer.from([i]))
      .digest();
    t.copy(okm, offset);
    offset += t.length;
  }
  return okm.slice(0, length);
}

/// Derive Track B directional key with multiplicity label.
export function deriveTrackBKey(
  sharedKey: Buffer,
  salt: Buffer,
  direction: Direction,
  ctxHex: string,
  multiplicityLabel: Buffer // 8 bytes
): Buffer {
  const info = Buffer.concat([
    Buffer.from("QSBS-v1:enc:"),
    Buffer.from(direction),
    Buffer.from(ctxHex.replace(/^0x/, ""), "hex"),
    multiplicityLabel,
  ]);
  const prk = hkdfExtract(salt, sharedKey);
  return hkdfExpand(prk, info, 32);
}

/// Compute Track B commitment handle (frequency + multiplicity-based salt).
export async function computeTrackBHandle(
  serializedFrequency: string,
  multiplicityIndex: number
): Promise<string> {
  await init(); // ensure WASM is ready
  const salt = `trackb-${multiplicityIndex.toString(16)}`;
  return trackb_commitment(serializedFrequency, salt);
}
\end{verbatim}

These snippets are illustrative and omit error handling, AEAD integration, and production hardening, but they show how Track B can be concretely wired into a BN254/WASM + TypeScript environment.

%-----------------------------------------------------------------------------%
\subsection{Implementation Guidelines and Recommendations}
%-----------------------------------------------------------------------------%

\subsubsection{Track B Tiers}

To make deployment practical, we recommend three implementation tiers for Track B:

\begin{description}[leftmargin=1.5em]
  \item[Tier 1: Commitment and receipts.] Implement BN254 Pedersen commitments for Track B states and publish \(h_t^{(B)}\) as audit receipts. No encryption is required at this tier.
  \item[Tier 2: Transcript and multiplicity binding.] Add per-role transcript chains, context hash, and multiplicity operator \(M_t^{(B)}\). Use them to derive AEAD keys and to label Track B events.
  \item[Tier 3: Full QSBS and feedback dynamics.] Integrate joint frequency mapping, coupling tensor \(T_t^{(B)}\), and the feedback map \(f_B\) controlling adaptive key rotation and convergence-based policy decisions.
\end{description}

\subsubsection{Fastest Validation Path}

A practical validation plan for Track B is:
\begin{enumerate}[leftmargin=1.5em]
  \item Implement the BN254 Pedersen commitment and a minimal Track B wrapper that computes \(h_t^{(B)}\) from serialized state.
  \item Design a small Track B protocol (e.g., three-stage state update) and record transcripts, multiplicity vectors, and commitments.
  \item Verify that commitments recompute correctly and that multiplicity indices match manually counted message-class occurrences.
  \item Add transcript chains and context hash; confirm that modifying any message changes both \(\mathrm{ctx}_t^{(B)}\) and the derived Track B keys.
  \item Introduce a simple feedback rule for key rotation (e.g., increase rotation frequency when latency or error rises) and empirically confirm the convergence properties in a simulation.
\end{enumerate}

%-----------------------------------------------------------------------------%
\subsection{Defensive Publication Statement}
%-----------------------------------------------------------------------------%

\subsubsection*{Scope and Intent}

This document constitutes a formal defensive publication covering Track B as defined above, including:

\begin{itemize}[leftmargin=1.5em]
  \item The use of BN254 Pedersen commitments over multiplicity-encoded frequency representations of Track B state.
  \item The embedding of prime-indexed multiplicity operator states into HKDF info strings for directional AEAD key derivation.
  \item The design of Track B as a protocol whose transcript history is G\"odel-encoded via prime exponents, enabling recursive stability and audit.
  \item The integration of a rank-1 coupling tensor and feedback-driven contraction map for adaptive Track B key rotation and convergence.
\end{itemize}

\subsubsection*{Public-Domain Defensive Release}

By publishing this description with sufficient detail for an expert to re-implement and analyze the described mechanisms, the authors intend to establish prior art and prevent third parties from obtaining patent rights that would block or encumber implementations of Track B conforming to this specification.
This includes, but is not limited to, patent claims over the specific combination of multiplicity operators, BN254 Pedersen commitments, transcript-bound HKDF derivation, and adaptive feedback dynamics as applied to certification and secure transport in Track B.

\section{Track C: WESAD CAS Proof-of-Concept \\
Defensive Prior Art Publication and Technical Report}


This section discloses a proof-of-concept application of the Codon-Aware Spectral (CAS) architecture to the WESAD wearable stress and affect detection dataset, constituting a Track~C extension of the Multiplicity Theory framework.\footnote{The WESAD dataset is a publicly available multimodal dataset for wearable stress and affect detection with 15 subjects and multiple physiological and motion channels, as described by Schmidt et al.\cite{Schmidt2018_WESAD,Schmidt2018_WESAD_UBI}}[web:25][web:408][web:405] The purpose of this publication is to establish prior art for the following contributions:

\begin{itemize}
  \item A CAS-style, channel-structured representation of WESAD signals using a fixed, low-dimensional axis system (Baseline, Stress, Amusement, and ancillary control channels).
  \item A Walsh--Hadamard-based spectral mixing layer on top of physiological features, treated as a CAS instance with orthogonal mixing and operator-norm guarantees.\cite{WHTOverview,WHTOrthogonal}[web:331][web:363][web:419]
  \item Integration of this spectral CAS representation with a Multiplicity-Crypto-compatible commitment and certification layer for stress/affect classification outputs (Track~B compatibility).
  \item A concrete Python reference implementation sketch demonstrating the full pipeline from raw WESAD signals to CAS features, classifier outputs, and commitment records, intended as a defense against subsequent patent claims over substantially similar architectures.\cite{WESAD_JohnsonGitHub,WESAD_MMIImplementation}[web:392][web:403]
\end{itemize}

The key novelty lies in treating WESAD as a CAS instance with (i) structured channels, (ii) orthogonal WHT-based spectral mixing, and (iii) optional attachment to multiplicity-crypto certification, thereby unifying stress-detection analytics with the broader Multiplicity Theory ecosystem.

\subsection{Background and Scope}

\subsubsection{WESAD Dataset}

The WESAD (Wearable Stress and Affect Detection) dataset comprises physiological and motion signals collected from 15 subjects via chest- and wrist-worn devices during lab sessions inducing neutral (baseline), stress, and amusement conditions.\cite{Schmidt2018_WESAD}[web:25][web:408][web:405] Available modalities include:

\begin{itemize}
  \item Electrocardiogram (ECG),
  \item Electrodermal activity (EDA),
  \item Electromyogram (EMG),
  \item Respiration (RESP),
  \item Skin temperature (TEMP),
  \item Blood volume pulse (BVP),
  \item Three-axis acceleration (ACC\_X, ACC\_Y, ACC\_Z).
\end{itemize}

The dataset has been used to benchmark stress vs.\ non-stress and multi-class affect detection (baseline vs.\ stress vs.\ amusement) using classical machine-learning and deep-learning models.\cite{Schmidt2018_WESAD,WESAD_JohnsonGitHub,WESAD_JaganGitHub,WESAD_Siyixu,WESAD_Srovins,WESAD_CrossModality,WESAD_CrossDomain}[web:408][web:399][web:393][web:395][web:396][web:22][web:418]

\subsubsection{CAS Architecture and WHT Layer}

In Track~A (Genomic CAS) and Track~B (Multiplicity-Crypto), CAS architectures used orthogonal Walsh--Hadamard-type transforms as mixing layers on top of structured channels.\cite{WHTOverview,WHTOrthogonal}[web:331][web:363][web:419] Track~C repurposes this machinery as follows:

\begin{enumerate}
  \item Define a small number of WESAD CAS channels capturing affective structure (baseline vs.\ stress vs.\ amusement), sensor-specific summary axes, and optional control axes (motion, artifacts).
  \item Map high-dimensional time-series features into these channels.
  \item Apply an orthogonal WHT mixing matrix to obtain a compact spectral representation.
  \item Use these CAS-spectral features as inputs to standard classifiers (e.g., logistic regression, SVM, shallow neural networks) for stress vs.\ non-stress and three-way affect detection.
  \item Optionally attach multiplicity-crypto commitments to model outputs and training events (Track~B integration).
\end{enumerate}

The remainder of this document states the CAS/WHT formulation for WESAD, states norm bounds, and provides a reference implementation sketch.

\subsection{Mathematical CAS Formulation for WESAD}

\subsubsection{Signal and Windowing Model}

Let \(s\) index subjects, \(c\) index sensor channels, and \(t\) denote discrete sample time.

\begin{definition}[Raw signal tensor]
Let
\[
  X_s(t) = \bigl( x_{s,c}(t) \bigr)_{c=1}^{C} \in \R^{C}
\]
be the multichannel time series for subject \(s\), where \(C\) aggregates all chosen WESAD modalities (e.g., ECG, EDA, RESP, TEMP, BVP, ACC components).\cite{Schmidt2018_WESAD,WESAD_CrossDomain}[web:408][web:418]
\end{definition}

We segment the signals into fixed-length windows.

\begin{definition}[Windowed segments]
Fix a window length \(L\) (in samples) and stride \(R\).
Define the \(j\)-th window for subject \(s\) as
\[
  X_{s,j} = \bigl(x_{s,c}(t)\bigr)_{c=1,\dots,C;\ t = t_j,\dots,t_j+L-1}
  \in \R^{C \times L},
\]
where \(t_j = t_0 + jR\).
Associate each window with a label \(y_{s,j} \in \{0,1,2\}\) for baseline, stress, and amusement, respectively.\cite{Schmidt2018_WESAD}[web:408][web:25]
\end{definition}

\subsubsection{Feature Extraction and Channel Aggregation}

For each sensor channel and window, we compute summary features such as mean, variance, spectral power, and domain-specific measures (e.g., HRV features from ECG, tonic/phasic components from EDA).\cite{WESAD_JohnsonGitHub,WESAD_MMIImplementation,WESAD_Siyixu}[web:392][web:403][web:395]

\begin{definition}[Feature vector]
For each window \((s,j)\), define a feature vector
\[
  f_{s,j} \in \R^d,
\]
built from concatenating features across all sensor modalities for that window.
\end{definition}

The CAS view is to group these features into a small number of channels.

\begin{definition}[Track C channel map]
Let \(\mathcal{K} = \{1,\dots,K\}\) index CAS channels.
Define disjoint index sets \(\mathcal{F}_k \subset \{1,\dots,d\}\) such that \(\bigcup_{k=1}^K \mathcal{F}_k = \{1,\dots,d\}\).
Define the channel aggregation map
\[
  A : \R^d \to \R^K, \quad
  (A f_{s,j})_k = \sum_{\ell \in \mathcal{F}_k} f_{s,j,\ell}.
\]
Then the channel vector for window \((s,j)\) is
\[
  a_{s,j} = A f_{s,j} \in \R^K.
\]
\end{definition}

Typical choices for channels include:

\begin{itemize}
  \item Cardio-respiratory channel (ECG, RESP-derived features),
  \item Electrodermal channel (EDA features),
  \item Motion channel (acceleration features),
  \item Temperature channel,
  \item Composite affective channels aligned with baseline, stress, and amusement axes (via supervised or unsupervised projection).
\end{itemize}

\subsubsection{Orthogonal WHT Mixing}

Let \(K = 2^m\) for some small \(m\), or let \(K\) be approximated by an orthogonal WHT-like matrix (as in Track~A) when not a power of two.\cite{WHTOverview,WHTOrthogonal}[web:331][web:363][web:419]

\begin{definition}[WESAD WHT mixing matrix]
Let \(H_K \in \R^{K \times K}\) be an orthogonal matrix (e.g., a normalized Walsh--Hadamard matrix or an orthogonal approximation when \(K\) is not a power of two).
Thus,
\[
  H_K H_K^\top = I_K = H_K^\top H_K.
\]
For each window \((s,j)\), define the CAS spectrum
\[
  h_{s,j} = H_K a_{s,j} \in \R^K.
\]
\end{definition}

\begin{proposition}[Norm preservation]
For all windows \((s,j)\),
\[
  \norm{h_{s,j}}_2 = \norm{a_{s,j}}_2.
\]
\end{proposition}

\begin{proof}
Because \(H_K\) is orthogonal, \(\norm{H_K x}_2 = \norm{x}_2\) for any \(x \in \R^K\).
This is standard for Walsh--Hadamard and other orthogonal transforms.\cite{WHTOverview,WHTOrthogonal}[web:331][web:363][web:419]
\end{proof}

\begin{proposition}[Aggregation operator norm bound]
If the aggregation matrix \(A\) satisfies \(\norm{A}_{2\to2} \le B_A\), then
\[
  \norm{h_{s,j}}_2 \le B_A \norm{f_{s,j}}_2.
\]
\end{proposition}

\begin{proof}
Since \(h_{s,j} = H_K A f_{s,j}\),
\[
  \norm{h_{s,j}}_2 \le \norm{H_K}_{2\to2} \norm{A}_{2\to2} \norm{f_{s,j}}_2 = B_A \norm{f_{s,j}}_2,
\]
because \(\norm{H_K}_{2\to2} = 1\) for an orthogonal matrix.
\end{proof}

Thus CAS-WHT features do not amplify norms beyond the aggregation operator norm, which is fixed given a channel design.

\subsection{Classification Layer and Certification Hooks}

\subsubsection{Feature Vectors for Stress/Affect Detection}

We consider one or more of the following feature families for stress/affect classification:

\begin{itemize}
  \item Raw aggregated channels \(a_{s,j}\),
  \item Spectral channels \(h_{s,j}\),
  \item Residuals or deviations from baseline (within-subject normalization),
  \item Concatenations of channel and spectral features.
\end{itemize}

\begin{definition}[CAS feature map]
Let \(\Phi : \R^d \to \R^D\) be the composition
\[
  \Phi(f_{s,j}) = \Psi\bigl( a_{s,j}, h_{s,j}\bigr),
\]
where \(\Psi\) is any fixed, differentiable map (e.g., stacking, residual computation, normalization).
The CAS feature for window \((s,j)\) is
\[
  z_{s,j} = \Phi(f_{s,j}) \in \R^D.
\]
\end{definition}

Because \(A\) and \(H_K\) are linear and bounded in operator norm, and \(\Psi\) can be designed to be Lipschitz, the overall map \(\Phi\) is Lipschitz, which simplifies stability and optimization analysis.

\subsubsection{Logistic Regression Model}

For binary stress vs.\ non-stress classification, we define:

\begin{definition}[Binary Track C classifier]
Let \(y_{s,j} \in \{0,1\}\) indicate non-stress vs.\ stress.
A logistic regression model uses parameters \(\theta \in \R^D\) and
\[
  p_{\theta}(y_{s,j}=1 \mid z_{s,j}) = \sigma(\theta^\top z_{s,j}),
\]
where \(\sigma(z) = 1/(1+e^{-z})\).
\end{definition}

The per-sample loss is
\[
  \ell(\theta; z_{s,j}, y_{s,j})
  = -y_{s,j} \log \sigma(\theta^\top z_{s,j})
    - (1-y_{s,j}) \log \bigl(1-\sigma(\theta^\top z_{s,j})\bigr).
\]

\begin{proposition}[Gradient bound]
If \(\norm{z_{s,j}}_2 \le B_Z\) for all \((s,j)\), then
\[
  \norm{\nabla_\theta \ell(\theta; z_{s,j}, y_{s,j})}_2 \le B_Z.
\]
\end{proposition}

\begin{proof}
The gradient is
\[
  \nabla_\theta \ell = \bigl(\sigma(\theta^\top z_{s,j}) - y_{s,j}\bigr) z_{s,j}.
\]
Because \(\sigma(\theta^\top z_{s,j}) \in (0,1)\) and \(y_{s,j} \in \{0,1\}\), we have \(\abs{\sigma(\theta^\top z_{s,j}) - y_{s,j}} \le 1\).
Thus
\[
  \norm{\nabla_\theta \ell}_2 \le \norm{z_{s,j}}_2 \le B_Z.
\]
\end{proof}

Using the previous norm bounds for \(A\) and \(H_K\), one can guarantee a fixed \(B_Z\) given bounded feature vectors \(f_{s,j}\).

\subsubsection{Certification Hooks (Track B Compatibility)}

To integrate with Track~B's multiplicity-crypto layer, Track~C can define:

\begin{itemize}
  \item A per-subject or per-session multiplicity operator \(M_t^{(C)}\) counting events such as model training epochs, threshold crossings, and classification outputs.
  \item Commitment values \(C(v,\rho)\) over CAS feature aggregates or model snapshots using Pedersen commitments in a BN254-like group.\cite{Pedersen_ZKDocs,Pedersen_RareSkills}[web:366][web:369]
  \item HKDF-based key derivation for encrypted logging of classification results, with WESAD CAS metadata included as associated data.
\end{itemize}

The detailed multiplicity-crypto machinery is beyond the scope of this appendix, but is structurally inherited from Track~B.

\subsection{Reference Implementation Sketch}

This section provides a Python-style code snippet for a minimal Track~C CAS pipeline on WESAD. Implementations in existing repositories demonstrate traditional pipelines, but do not use a CAS/WHT structure.\cite{WESAD_JohnsonGitHub,WESAD_MMIImplementation}[web:392][web:403]

\subsubsection{Data Loading and Preprocessing}

\begin{lstlisting}[language=Python,caption={Minimal WESAD CAS feature pipeline (conceptual).}]
import numpy as np
from scipy.signal import welch
from scipy.linalg import hadamard

# --- Configuration ---

WINDOW_SECONDS = 60
STRIDE_SECONDS = 30
FS_ECG = 700.0  # example sampling rate, see WESAD documentation
K_CHANNELS = 8  # number of CAS channels (power of 2 for standard WHT)

# --- Helper: windowing ---

def sliding_windows(x, window, stride):
    """Simple sliding-window generator along axis=0."""
    start = 0
    n = x.shape[0]
    while start + window <= n:
        yield x[start:start + window]
        start += stride

# --- Feature extraction per window ---

def extract_features_window(ecg, eda, resp, temp, acc):
    """Return feature vector f in R^d for one window."""
    feats = []

    # Example ECG features (time-domain + frequency-domain)
    feats.append(np.mean(ecg))
    feats.append(np.std(ecg))
    f_ecg, p_ecg = welch(ecg, fs=FS_ECG, nperseg=len(ecg))
    feats.append(np.trapz(p_ecg[(f_ecg >= 0.04) & (f_ecg <= 0.15)]) )  # LF power
    feats.append(np.trapz(p_ecg[(f_ecg >= 0.15) & (f_ecg <= 0.4)]) )   # HF power

    # EDA features
    feats.append(np.mean(eda))
    feats.append(np.std(eda))

    # Respiration features
    feats.append(np.mean(resp))
    feats.append(np.std(resp))

    # Temperature features
    feats.append(np.mean(temp))

    # ACC features (3-axis)
    for axis in range(acc.shape[1]):
        feats.append(np.mean(acc[:, axis]))
        feats.append(np.std(acc[:, axis]))

    return np.array(feats, dtype=np.float32)  # f in R^d

# --- Channel aggregation ---

def build_channel_matrix(d, K):
    """
    Build aggregation matrix A in R^{K x d}.
    Here we use a simple block assignment; in practice, this
    can be hand-crafted to group semantically related features.
    """
    A = np.zeros((K, d), dtype=np.float32)
    block = int(np.ceil(d / K))
    idx = 0
    for k in range(K):
        for j in range(block):
            if idx < d:
                A[k, idx] = 1.0
                idx += 1
    return A

# --- WHT mixing ---

def hadamard_mixing(K):
    """
    Return orthogonal mixing matrix H_K in R^{K x K}.
    For K a power of 2, use normalized Hadamard matrix.
    """
    H = hadamard(K).astype(np.float32)
    H /= np.sqrt(K)
    return H

# --- Main CAS feature pipeline for one subject ---

def cas_features_for_subject(subj_data):
    """
    subj_data: dict with keys 'ecg', 'eda', 'resp', 'temp', 'acc', 'labels'
      Each signal is a 1D or 2D numpy array aligned in time.
    Returns:
      Z: CAS features in R^{num_windows x D}
      Y: integer labels per window
    """
    ecg = subj_data["ecg"]
    eda = subj_data["eda"]
    resp = subj_data["resp"]
    temp = subj_data["temp"]
    acc = subj_data["acc"]  # shape: (T, 3)
    labels = subj_data["labels"]  # per-sample label

    win = int(WINDOW_SECONDS * FS_ECG)
    stride = int(STRIDE_SECONDS * FS_ECG)

    # Build A and H_K once
    # First, get d from a dummy window
    f0 = extract_features_window(ecg[:win], eda[:win], resp[:win], temp[:win], acc[:win])
    d = f0.shape[0]
    A = build_channel_matrix(d, K_CHANNELS)
    H = hadamard_mixing(K_CHANNELS)

    Z = []
    Y = []

    for w_idx, (ecg_w, eda_w, resp_w, temp_w, acc_w) in enumerate(
        zip(sliding_windows(ecg, win, stride),
            sliding_windows(eda, win, stride),
            sliding_windows(resp, win, stride),
            sliding_windows(temp, win, stride),
            sliding_windows(acc, win, stride))
    ):
        f = extract_features_window(ecg_w, eda_w, resp_w, temp_w, acc_w)   # in R^d
        a = A @ f                                                           # in R^K
        h = H @ a                                                           # in R^K

        # Simple CAS feature: concatenate a and h
        z = np.concatenate([a, h], axis=0)
        Z.append(z)

        # Majority label within window
        label_w = np.bincount(labels[w_idx * stride: w_idx * stride + win]).argmax()
        Y.append(label_w)

    return np.stack(Z, axis=0), np.array(Y, dtype=np.int64)
\end{lstlisting}

This code demonstrates:

\begin{itemize}
  \item Construction of windowed WESAD features from ECG, EDA, RESP, TEMP, and ACC.\cite{WESAD_JohnsonGitHub,WESAD_MMIImplementation}[web:392][web:403]
  \item A channel aggregation matrix \(A\) and orthogonal WHT mixing \(H\) as in the CAS formalism.
  \item A final CAS feature vector per window \(z_{s,j} \in \R^{2K}\).
\end{itemize}

\subsubsection{Classifier and Certification Hook}

\begin{lstlisting}[language=Python,caption={Classifier training and placeholder for certification.}]
from sklearn.linear_model import LogisticRegression
from sklearn.model_selection import train_test_split
from sklearn.metrics import classification_report

# Suppose we aggregate all subjects:
Z_all, Y_all = [], []
for subj_data in all_subjects_data:  # list of subject dicts
    Z_subj, Y_subj = cas_features_for_subject(subj_data)
    Z_all.append(Z_subj)
    Y_all.append(Y_subj)
Z_all = np.concatenate(Z_all, axis=0)
Y_all = np.concatenate(Y_all, axis=0)

# Binary label example: stress vs. non-stress
Y_bin = (Y_all == STRESS_LABEL).astype(int)

X_train, X_test, y_train, y_test = train_test_split(
    Z_all, Y_bin, test_size=0.2, stratify=Y_bin, random_state=42
)

clf = LogisticRegression(max_iter=1000)
clf.fit(X_train, y_train)

y_pred = clf.predict(X_test)
print(classification_report(y_test, y_pred))

# --- Certification hook (conceptual) ---

def pedersen_commit(value_bytes, blinding_scalar):
    # placeholder stub for BN254 Pedersen commitment
    # using external library in a real implementation
    # C = rho * G + v * H
    pass

# Example: commit to model parameters
theta_bytes = clf.coef_.astype(np.float32).tobytes() + clf.intercept_.astype(np.float32).tobytes()
rho = os.urandom(32)  # real implementation: map to scalar mod r
model_commitment = pedersen_commit(theta_bytes, rho)
\end{lstlisting}

In a full Track~C system, the model commitment would be bound to training transcript and evaluation results using the Track~B multiplicity-crypto scheme (HKDF-derived keys, multiplicity operator, and transcript chains).

\subsection{Prior Art Scope and Claims}

This document intends to place the following objects and constructions into the public domain as prior art:

\begin{enumerate}
  \item \textbf{CAS representation for WESAD}: The notion of defining a finite set of WESAD CAS channels, aggregating high-dimensional time-series features into these channels, and applying an orthogonal Walsh--Hadamard or WHT-like transform to obtain a low-dimensional spectral representation for stress/affect detection.\cite{Schmidt2018_WESAD,WHTOverview,WHTOrthogonal}[web:408][web:331][web:363][web:419]
  \item \textbf{Operator-norm guarantees}: The use of explicit operator-norm bounds on the channel aggregation matrix and orthogonal transform to guarantee stability and boundedness of CAS features when used in logistic or related classifiers.
  \item \textbf{Track~B-compatible certification hooks}: The explicit inclusion of multiplicity-counting operators and Pedersen-style commitments on top of WESAD CAS features and model parameters, as a pattern for certified stress/affect analytics, even when not fully implemented in this proof-of-concept.\cite{Pedersen_ZKDocs,Pedersen_RareSkills}[web:366][web:369]
  \item \textbf{Reference implementation}: The presented Python pipeline structure from raw WESAD data to CAS features, classifier training, and commitment hooks, as a specific algorithmic template.
\end{enumerate}

These items, together with the mathematical derivations and code sketches given above, should be considered prior art against any attempt to claim exclusive rights over substantially similar combinations of WESAD, CAS, WHT, and multiplicity-crypto certification.

\appendix
\section{Track A: Hadamard Matrices and 6-Channel WHT Approximation}

\subsection{Classical Hadamard Matrices}

The Walsh--Hadamard transform is defined by a family of \emph{Hadamard matrices}, which are real orthogonal matrices with entries of equal magnitude.\cite{HadamardOverview,HadamardOrthonormal}

\begin{definition}[Hadamard matrix]
A (Sylvester-type) Hadamard matrix of order \(n = 2^m\) is a matrix \(H_n \in \{-1,1\}^{n \times n}\) defined recursively by
\[
  H_1 = [1], \qquad
  H_{2n} = \begin{bmatrix}
    H_n & H_n \\[0.3em]
    H_n & -H_n
  \end{bmatrix}.
\]
Its normalized version is
\[
  \widehat{H}_n = \frac{1}{\sqrt{n}} H_n \in \R^{n \times n}.
\]
\end{definition}

\begin{proposition}[Orthogonality and norm preservation]
For each \(n = 2^m\),
\[
  \widehat{H}_n \widehat{H}_n^\top = I_n = \widehat{H}_n^\top \widehat{H}_n,
\]
so \(\widehat{H}_n\) is orthonormal. In particular, for any \(x \in \R^n\),
\[
  \norm{\widehat{H}_n x}_2 = \norm{x}_2.
\]
\end{proposition}

\begin{proof}
The orthogonality property follows by induction on \(n\), using the block recursion and the fact that for the base case \(H_1 = [1]\),
\(\widehat{H}_1 = [1]\) is trivially orthonormal. Standard references on the Walsh--Hadamard transform derive these identities explicitly and highlight that the normalized matrix has entries of magnitude \(n^{-1/2}\) and is orthonormal.\cite{HadamardOverview,HadamardOrthonormal} The 2-norm preservation is a direct consequence of orthonormality.
\end{proof}

\subsubsection{Restriction to Six Channels}

Track A does not operate in a power-of-two dimension; instead, it aggregates 61 codon frequencies into six biologically motivated channels and then applies a \emph{6-channel WHT-like} transformation.

\begin{definition}[6-channel transform]
A matrix \(H_6 \in \R^{6 \times 6}\) is called a \emph{6-channel WHT approximation} if:
\begin{enumerate}
  \item \(H_6\) is orthogonal: \(H_6 H_6^\top = I_6 = H_6^\top H_6\).
  \item Entries of \(H_6\) have equal magnitude up to a common scaling (e.g., \(\pm 1/\sqrt{6}\)), mirroring Hadamard structure.
\end{enumerate}
Given a channel intensity vector \(a_g \in \R^6\) for gene \(g\), the 6-channel spectrum is
\[
  h_g = H_6 a_g \in \R^6.
\]
\end{definition}

\begin{proposition}[Norm preservation for 6-channel transform]
If \(H_6\) is orthogonal, then for each gene \(g\),
\[
  \norm{h_g}_2 = \norm{a_g}_2.
\]
\end{proposition}

\begin{proof}
Orthogonality gives \(H_6^\top H_6 = I_6\). For any \(a_g\),
\[
  \norm{h_g}_2^2 = h_g^\top h_g = a_g^\top H_6^\top H_6 a_g = a_g^\top a_g = \norm{a_g}_2^2.
\]
Taking square roots yields the claim.
\end{proof}

\begin{remark}
Although a true Hadamard matrix exists only for dimensions \(2^m\), any real orthogonal \(6 \times 6\) matrix can serve as the spectral mixing matrix for Track A, and its orthogonality is all that is required for the norm bounds and distance properties used in the main text.\cite{HadamardOverview,HadamardOrthonormal}
\end{remark}

\subsection{Codon-Usage Vectors and Channel Aggregation}

\subsubsection{Codon-Usage Geometry}

Let \(\mathcal{C}\) be the set of 61 sense codons. For each gene \(g\),
the normalized codon-usage vector \(c_g \in \R^{61}\) satisfies
\[
  c_g(c) \ge 0 \quad \forall c \in \mathcal{C},
  \qquad \sum_{c \in \mathcal{C}} c_g(c) = 1,
\]
so \(c_g\) lies in the 61-dimensional simplex.

We partition \(\mathcal{C}\) into six disjoint channel sets \(\mathcal{C}_1, \dots, \mathcal{C}_6\), chosen to reflect biologically meaningful groupings (GC3, amino-acid classes, etc.) as described in the main specification.

\begin{definition}[Channel aggregation map]
The channel aggregation map \(A : \R^{61} \to \R^6\) is defined by
\[
  a_{g,i} = \sum_{c \in \mathcal{C}_i} c_g(c), \quad i = 1,\dots,6,
\]
and \(a_g = (a_{g,1},\dots,a_{g,6})^\top\).
\end{definition}

\begin{proposition}[Norm bounds under aggregation]
Let \(\norm{\cdot}_2\) be the Euclidean norm on both \(\R^{61}\) and \(\R^6\).
Then for each gene \(g\),
\[
  \norm{a_g}_2 \le \norm{c_g}_2 \le 1.
\]
\end{proposition}

\begin{proof}
Denote the aggregation matrix by \(A \in \{0,1\}^{6 \times 61}\), where
\((A)_{i,c} = 1\) if \(c \in \mathcal{C}_i\) and \(0\) otherwise.
Then \(a_g = A c_g\).

Each row of \(A\) has disjoint support and entries in \(\{0,1\}\), so
\[
  (A A^\top)_{i,i} = \sum_{c} (A_{i,c})^2 = |\mathcal{C}_i|,
  \quad
  (A A^\top)_{i,j} = 0 \quad (i \ne j).
\]
Thus \(A A^\top = \mathrm{diag}(|\mathcal{C}_1|,\dots,|\mathcal{C}_6|)\), and the operator norm of \(A\) satisfies
\[
  \norm{A}_{2\to2}^2 = \lambda_{\max}(A A^\top) = \max_i |\mathcal{C}_i|.
\]
In practice each \(|\mathcal{C}_i|\) is finite and modest, and since
\(\sum_{c} c_g(c) = 1\) with all components nonnegative, we have
\(\norm{c_g}_2 \le \sqrt{\sum_c c_g(c)^2} \le \sqrt{\sum_c c_g(c)} = 1\). Hence
\[
  \norm{a_g}_2 = \norm{A c_g}_2 \le \norm{A}_{2\to2}\norm{c_g}_2 \le \sqrt{\max_i |\mathcal{C}_i|}.
\]
Because \(\max_i |\mathcal{C}_i|\) is a constant (independent of \(g\)), we can absorb this into the global constant of the Track A operator-norm budget.
For scaling convenience, one can simply renormalize \(A\) to ensure \(\norm{A}_{2\to2} \le 1\), which yields \(\norm{a_g}_2 \le \norm{c_g}_2 \le 1\).
\end{proof}

\subsubsection{Composition with 6-Channel WHT}

Recall \(h_g = H_6 a_g\), where \(H_6\) is orthogonal.

\begin{proposition}[Overall spectral norm bound]
For each gene \(g\),
\[
  \norm{h_g}_2 \le \norm{c_g}_2 \le 1.
\]
\end{proposition}

\begin{proof}
Using the previous bounds and orthogonality of \(H_6\),
\[
  \norm{h_g}_2 = \norm{H_6 a_g}_2 \le \norm{H_6}_{2\to2} \norm{a_g}_2 = \norm{a_g}_2,
\]
since \(\norm{H_6}_{2\to2} = 1\) for an orthogonal matrix.
Taking \(\norm{a_g}_2 \le \norm{c_g}_2 \le 1\) from the previous proposition yields the claim.
\end{proof}

\subsection{GC3-Controlled Null Ensemble and \texorpdfstring{$\Delta h$}{Δh}}

\subsubsection{Null Ensemble Construction}

For each gene \(g\), the GC3-controlled null ensemble consists of random codon-usage realizations preserving the amino-acid sequence and GC3 content.

\begin{definition}[Null ensemble]
Fix a number of replicates \(R \in \N\). For gene \(g\), let \(\{c_{g,r}^{\mathrm{null}}\}_{r=1}^R\) be an i.i.d.\ sample of codon-usage vectors constructed by GC3-preserving codon reshuffling per position.
Define the corresponding aggregated spectra
\[
  h_{g,r}^{\mathrm{null}} = H_6 A c_{g,r}^{\mathrm{null}}, \quad r=1,\dots,R.
\]
\end{definition}

\begin{definition}[Mean null spectrum and deviation]
The mean null spectrum is
\[
  \bar{h}_g^{\mathrm{null}} = \frac{1}{R} \sum_{r=1}^R h_{g,r}^{\mathrm{null}}.
\]
Define the deviation vector
\[
  \Delta h_g = h_g - \bar{h}_g^{\mathrm{null}} \in \R^6,
\]
and its scalar magnitude
\[
  \delta_g = \norm{\Delta h_g}_2.
\]
\end{definition}

\subsubsection{Norm Bounds for \texorpdfstring{$\Delta h$}{Δh}}

\begin{proposition}[Deviation norm bound]
For each gene \(g\),
\[
  \norm{\Delta h_g}_2 \le \norm{h_g}_2 + \norm{\bar{h}_g^{\mathrm{null}}}_2 \le 2.
\]
\end{proposition}

\begin{proof}
By the triangle inequality,
\[
  \norm{\Delta h_g}_2 = \norm{h_g - \bar{h}_g^{\mathrm{null}}}_2 \le \norm{h_g}_2 + \norm{\bar{h}_g^{\mathrm{null}}}_2.
\]
Each null spectrum \(h_{g,r}^{\mathrm{null}}\) has the same norm bound as \(h_g\) because GC3-preserving reshuffles are still normalized codon-usage distributions; hence \(\norm{h_{g,r}^{\mathrm{null}}}_2 \le 1\) using the previous section's bound.
Therefore,
\[
  \norm{\bar{h}_g^{\mathrm{null}}}_2 =
  \left\lVert \frac{1}{R} \sum_{r=1}^R h_{g,r}^{\mathrm{null}} \right\rVert_2
  \le \frac{1}{R} \sum_{r=1}^R \norm{h_{g,r}^{\mathrm{null}}}_2
  \le 1.
\]
Combining with \(\norm{h_g}_2 \le 1\) yields \(\norm{\Delta h_g}_2 \le 2\).
\end{proof}

\begin{corollary}[Bound on scalar \texorpdfstring{$\delta_g$}{δg}]
For each gene \(g\), the scalar deviation satisfies
\[
  0 \le \delta_g \le 2.
\]
\end{corollary}

\begin{proof}
By definition \(\delta_g = \norm{\Delta h_g}_2 \ge 0\).
The upper bound follows directly from the previous proposition.
\end{proof}

\begin{remark}
In practice, because both \(h_g\) and \(\bar{h}_g^{\mathrm{null}}\) are weighted combinations of channels constructed from a probability vector, typical norms will be substantially below 1; the bound above is worst-case and sufficient for numerical stability and feature scaling guarantees in the Track A logistic models.
\end{remark}

\subsection{Logistic Regression and Feature Scaling}

\subsubsection{Feature Space and Scaling}

Let \(y_g \in \{0,1\}\) be the binary label for gene \(g\), and consider two feature sets:
\[
  X_g^{\mathrm{GC3}} = (\mathrm{GC3}_g) \in \R,
  \quad
  X_g^{\mathrm{GC3}+\delta} = (\mathrm{GC3}_g, \delta_g) \in \R^2,
\]
or the vector form
\[
  X_g^{\mathrm{GC3}+\Delta h} = (\mathrm{GC3}_g, \Delta h_{g,1}, \dots, \Delta h_{g,6}) \in \R^7.
\]

From the definition of GC3, we have \(0 \le \mathrm{GC3}_g \le 1\) for all genes.
Combined with \(0 \le \delta_g \le 2\), all coordinates of Track A feature vectors lie in a compact interval.
This uniform boundedness is useful for controlling the Lipschitz constants of logistic losses and for regularization tuning.

\begin{proposition}[Uniform boundedness of features]
For all genes \(g\),
\[
  \norm{X_g^{\mathrm{GC3}}}_2 \le 1,
  \quad
  \norm{X_g^{\mathrm{GC3}+\delta}}_2 \le \sqrt{1^2 + 2^2} = \sqrt{5},
\]
and similarly \(\norm{X_g^{\mathrm{GC3}+\Delta h}}_2 \le \sqrt{1 + 6 \cdot 2^2} = \sqrt{25} = 5\).
\end{proposition}

\begin{proof}
Immediate from the coordinate bounds: \(\mathrm{GC3}_g \in [0,1]\), each \(\Delta h_{g,i}\) component has magnitude at most \(\norm{\Delta h_g}_2 \le 2\), and \(\delta_g \le 2\).
\end{proof}

\subsubsection{Logistic Loss Lipschitz Bound}

Consider the logistic regression model with parameter vector \(\theta \in \R^d\) and features \(X_g \in \R^d\),
\[
  p_\theta(y_g=1 \mid X_g) = \sigma(\theta^\top X_g), \quad
  \sigma(z) = \frac{1}{1 + e^{-z}}.
\]

Define the negative log-likelihood for a single sample:
\[
  \ell(\theta; X_g, y_g) = -y_g \log \sigma(\theta^\top X_g) - (1-y_g)\log (1-\sigma(\theta^\top X_g)).
\]

\begin{proposition}[Gradient Lipschitz constant]
Let \(X_g\) be any of the Track A feature vectors described above, with \(\norm{X_g}_2 \le B\) for all \(g\).
Then the gradient of \(\ell\) with respect to \(\theta\) satisfies
\[
  \norm{\nabla_\theta \ell(\theta; X_g, y_g)}_2 \le B.
\]
\end{proposition}

\begin{proof}
The gradient is
\[
  \nabla_\theta \ell(\theta; X_g, y_g) = \bigl(\sigma(\theta^\top X_g) - y_g\bigr) X_g.
\]
Since \(\sigma(\theta^\top X_g) \in (0,1)\) and \(y_g \in \{0,1\}\), we have \(\abs{\sigma(\theta^\top X_g) - y_g} \le 1\).
Therefore
\[
  \norm{\nabla_\theta \ell(\theta; X_g, y_g)}_2
  = \abs{\sigma(\theta^\top X_g) - y_g}\,\norm{X_g}_2
  \le \norm{X_g}_2 \le B.
\]
\end{proof}

\begin{corollary}[Global Lipschitz bound]
Over the Track A feature sets, the logistic loss gradient is uniformly bounded with \(B = 5\) in the worst case (GC3+\(\Delta h\) vector model), which yields global step-size conditions for first-order optimization methods.
\end{corollary}

\subsection{Contraction Arguments for Feedback Variants}

While Track A in its simplest form does not require dynamic feedback, the same mathematical framework used in Track B can be restricted to the purely classical Track A setting, should adaptive weighting or iterative re-labeling be introduced.

\subsubsection{Static vs.\ Dynamic Interpretation}

In a dynamic interpretation, one might define a Track A state \(S_t^{(A)} = (c_{g,t}, h_{g,t}, \Delta h_{g,t})\) evolving under iterative refinement (e.g., repeated null-ensemble updates or re-weighting of channels based on current classifier residuals).

\begin{definition}[Track A feedback map (hypothetical)]
Let \(Z_t^{(A)}\) denote a vector of Track A state descriptors (e.g., averaged spectra, residuals, etc.).
A feedback map is a function
\[
  f_A : \R^d \times \R^k \to \R^d, \quad Z_{t+1}^{(A)} = f_A(Z_t^{(A)}, \Xi_t),
\]
where \(\Xi_t\) encodes observed statistics (e.g., accuracy, cross-entropy, etc.) at iteration \(t\).
\end{definition}

\begin{theorem}[Contraction condition for Track A feedback]
Suppose there exists \(L < 1\) such that
\[
  \norm{f_A(Z, \Xi) - f_A(Z', \Xi)}_2 \le L \norm{Z - Z'}_2
\]
for all \(Z,Z'\) in the state space and for all admissible \(\Xi\).
Then:
\begin{enumerate}
  \item There exists a unique fixed point \(Z_\ast^{(A)}\).
  \item For any initialization \(Z_0^{(A)}\), the iterates converge geometrically:
  \[
    \norm{Z_t^{(A)} - Z_\ast^{(A)}}_2 \le L^t \norm{Z_0^{(A)} - Z_\ast^{(A)}}_2.
  \]
\end{enumerate}
\end{theorem}

\begin{proof}
This is a direct application of the Banach Fixed-Point Theorem on the complete metric space \((\R^d, \norm{\cdot}_2)\).
\end{proof}

Although Track A as specified is a single-pass pipeline, this result is included to document the operator-norm and contraction logic underlying any future adaptive extensions of the protocol.

\subsection{Summary of Operator Norm Bounds}

For convenience, we collect the key norm bounds proved above.

\begin{itemize}
  \item \textbf{Hadamard/WHT norm preservation.}
  For normalized Hadamard matrices \(\widehat{H}_n\) and the 6-channel orthogonal matrix \(H_6\),
  \[
    \norm{\widehat{H}_n x}_2 = \norm{x}_2, \quad
    \norm{H_6 a}_2 = \norm{a}_2.
  \]
  \item \textbf{Codon aggregation bound.}
  With appropriate normalization of the aggregation matrix \(A\),
  \[
    \norm{a_g}_2 \le \norm{c_g}_2 \le 1.
  \]
  \item \textbf{6-channel spectrum bound.}
  For each gene,
  \[
    \norm{h_g}_2 \le \norm{c_g}_2 \le 1.
  \]
  \item \textbf{Null ensemble deviation.}
  For each gene,
  \[
    0 \le \delta_g = \norm{h_g - \bar{h}_g^{\mathrm{null}}}_2 \le 2.
  \]
  \item \textbf{Feature vector bounds.}
  For GC3, GC3+\(\delta\), and GC3+\(\Delta h\) feature vectors,
  \[
    \norm{X_g^{\mathrm{GC3}}}_2 \le 1, \quad
    \norm{X_g^{\mathrm{GC3}+\delta}}_2 \le \sqrt{5}, \quad
    \norm{X_g^{\mathrm{GC3}+\Delta h}}_2 \le 5.
  \]
  \item \textbf{Logistic loss gradient.}
  The gradient of the per-sample logistic loss satisfies
  \[
    \norm{\nabla_\theta \ell(\theta; X_g, y_g)}_2 \le B,
  \]
  with \(B \in \{1,\sqrt{5},5\}\) depending on the feature set.
\end{itemize}

These results provide the mathematical backbone for Track A’s numerical stability, scaling behavior, and potential extension to adaptive or iterative schemes.

\vspace{1em}
\noindent\textbf{References (for background only)}
\begin{itemize}
  \item Standard descriptions of Hadamard/Walsh--Hadamard transforms and their orthonormality properties.\cite{HadamardOverview,HadamardOrthonormal,HadamardMatrixTopic}
  \item Overviews of codon usage analysis and GC3-driven models, highlighting the central role of GC3 in codon usage variation.\cite{GCUA,GC3Model,GC3Variation}
\end{itemize}


\section{Track B: Groups, Commitments, and Basic Bounds}

\subsection{BN254 and Pedersen Commitments}

Let \(\G_1\) be the prime-order subgroup of an elliptic curve (e.g., BN254) with prime order \(r\) and generator \(G\).
Let \(\F_r\) denote the associated scalar field.
We assume the standard hardness of the discrete logarithm problem in \(\G_1\).

\begin{definition}[Independent generator \(H\)]
Let \(\mathsf{HashToScalar} : \{0,1\}^* \to \F_r\) be a random-oracle hash.
Define
\[
  h_G = \mathsf{HashToScalar}(\texttt{"pedersen\_H"} \,\|\, \mathrm{enc}(G)) \in \F_r,
\]
and set
\[
  H = h_G \cdot G \in \G_1.
\]
\end{definition}

In the random-oracle model and under the discrete logarithm assumption, this gives a generator \(H\) that is computationally independent of \(G\).\cite{PedersenZKDocs,PedersenRareSkills,CommitmentSchemeWiki}

\begin{definition}[BN254 Pedersen commitment]
For value \(v \in \F_r\) and blinding factor \(\rho \in \F_r\), the Pedersen commitment is
\[
  C(v,\rho) = \rho G + v H \in \G_1.
\]
\end{definition}

\begin{proposition}[Additive homomorphism]
For any \(v_1,v_2,\rho_1,\rho_2 \in \F_r\),
\[
  C(v_1,\rho_1) + C(v_2,\rho_2) = C(v_1+v_2,\rho_1+\rho_2).
\]
\end{proposition}

\begin{proof}
Direct expansion:
\[
  C(v_1,\rho_1) + C(v_2,\rho_2) = (\rho_1 G + v_1 H) + (\rho_2 G + v_2 H)
  = (\rho_1+\rho_2)G + (v_1+v_2)H = C(v_1+v_2,\rho_1+\rho_2).
\]
This is the standard additive homomorphism of Pedersen commitments.\cite{PedersenZKDocs,PedersenRareSkills}
\end{proof}

\begin{proposition}[Hiding and binding (informal)]
Assuming the discrete logarithm problem in \(\G_1\) is hard and \(\rho\) is sampled uniformly at random from \(\F_r\), the scheme \(C(v,\rho)\) is:
\begin{itemize}
  \item Computationally \emph{hiding}: given \(C(v,\rho)\), distinguishing \(v\) from any other \(v'\) is infeasible.
  \item Computationally \emph{binding}: it is infeasible to find \((v,\rho) \ne (v',\rho')\) such that \(C(v,\rho) = C(v',\rho')\).
\end{itemize}
\end{proposition}

\begin{proof}[Proof sketch]
These properties are standard for Pedersen commitments in any group of prime order with independent generators.\cite{PedersenZKDocs,PedersenRareSkills,PedersenLibby,PedersenLibbitcoin}
Binding reduces to the hardness of discrete logarithm in \(\G_1\); hiding follows from the fact that \(\rho G\) is uniformly distributed in \(\G_1\) when \(\rho\) is uniform in \(\F_r\).
\end{proof}

\subsection{Multiplicity Operator and Transcript Encoding}

\subsubsection{Prime-Indexed Encoding}

Let \(\{p_i\}_{i \ge 1}\) be the increasing sequence of prime numbers.
Let \(\mathcal{C} = \{c_1,\dots,c_k\}\) denote the finite set of Track B message classes (e.g., \texttt{state\_update}, \texttt{ack}, \texttt{key\_rotate}, etc.).

\begin{definition}[Class counts]
Given a Track B transcript \(\sigma_t^{(B)}\) up to time \(t\), define
\[
  c_i(\sigma_t^{(B)}) \in \N
\]
as the number of messages of class \(c_i\) that appear in \(\sigma_t^{(B)}\).
\end{definition}

\begin{definition}[Prime-indexed transcript encoding]
The transcript encoding at time \(t\) is
\[
  \mathrm{enc}_B(\sigma_t^{(B)}) = \prod_{i=1}^k p_i^{c_i(\sigma_t^{(B)})} \in \N_{>0}.
\]
\end{definition}

\begin{definition}[Track B multiplicity operator]
The Track B multiplicity operator is the vector
\[
  M_t^{(B)} = \bigl(c_1(\sigma_t^{(B)}), \dots, c_k(\sigma_t^{(B)})\bigr) \in \N^k.
\]
\end{definition}

\begin{proposition}[Injectivity]
If \(\sigma_t^{(B)}\) and \(\sigma_t'{}^{(B)}\) have different multisets of message classes, then
\(\mathrm{enc}_B(\sigma_t^{(B)}) \ne \mathrm{enc}_B(\sigma_t'{}^{(B)})\).
Equivalently, \(\mathrm{enc}_B\) is injective with respect to the message-class multiset.
\end{proposition}

\begin{proof}
By the Fundamental Theorem of Arithmetic, the prime factorization of a positive integer is unique.
If two encodings were equal, the exponent of each prime \(p_i\) in the factorization would have to match, so \(c_i(\sigma_t^{(B)}) = c_i(\sigma_t'{}^{(B)})\) for all \(i\).
Thus the class multiplicities are equal and the multisets coincide.
\end{proof}

\subsubsection{Norm Bounds for \texorpdfstring{$M_t^{(B)}$}{M\_t\^B}}

Let
\[
  N_t = \sum_{i=1}^k c_i(\sigma_t^{(B)}) = \text{total number of messages in the transcript up to time } t.
\]

\begin{proposition}[Norm bounds on \(M_t^{(B)}\)]
For each time \(t\),
\begin{align*}
  \norm{M_t^{(B)}}_1 &= N_t, \\
  \norm{M_t^{(B)}}_2 &\le N_t, \\
  \norm{M_t^{(B)}}_\infty &\le N_t.
\end{align*}
\end{proposition}

\begin{proof}
The \(\ell_1\) norm is
\[
  \norm{M_t^{(B)}}_1 = \sum_{i=1}^k \abs{c_i(\sigma_t^{(B)})} = \sum_{i=1}^k c_i(\sigma_t^{(B)}) = N_t.
\]
For the \(\ell_\infty\) norm,
\[
  \norm{M_t^{(B)}}_\infty = \max_i c_i(\sigma_t^{(B)}) \le \sum_i c_i(\sigma_t^{(B)}) = N_t.
\]
For the \(\ell_2\) norm,
\[
  \norm{M_t^{(B)}}_2^2 = \sum_{i=1}^k c_i(\sigma_t^{(B)})^2 \le \max_i c_i(\sigma_t^{(B)}) \sum_i c_i(\sigma_t^{(B)}) \le N_t^2,
\]
so \(\norm{M_t^{(B)}}_2 \le N_t\).
\end{proof}

When used as a diagonal linear operator (see below), these bounds transfer directly to operator norms.

\subsubsection{Multiplicity as a Linear Operator}

To use \(M_t^{(B)}\) in operator-norm bounds, we view it as a diagonal matrix.

\begin{definition}[Diagonal multiplicity operator]
Given \(M_t^{(B)}\), define the diagonal matrix
\[
  \mathbf{M}_t^{(B)} = \mathrm{diag}\bigl(M_t^{(B)}\bigr) \in \R^{k \times k}.
\]
\end{definition}

\begin{proposition}[Operator norms of \(\mathbf{M}_t^{(B)}\)]
The spectral norm (operator 2-norm) of \(\mathbf{M}_t^{(B)}\) is
\[
  \norm{\mathbf{M}_t^{(B)}}_{2\to2} = \norm{M_t^{(B)}}_\infty \le N_t,
\]
and more generally \(\norm{\mathbf{M}_t^{(B)}}_{p\to p} = \norm{M_t^{(B)}}_\infty\) for any \(p \in [1,\infty]\).
\end{proposition}

\begin{proof}
For a diagonal matrix \(D = \mathrm{diag}(d_1,\dots,d_k)\), the operator norm \(\norm{D}_{p\to p}\) is \(\max_i \abs{d_i}\) for any \(p\).
This is immediate by considering unit vectors aligned with coordinates.
Thus
\[
  \norm{\mathbf{M}_t^{(B)}}_{2\to2} = \max_i \abs{c_i(\sigma_t^{(B)})} = \norm{M_t^{(B)}}_\infty \le N_t.
\]
\end{proof}

\subsection{Coupling Tensor and Encryption Operator}

\subsubsection{Rank-1 Coupling Tensor}

In the QSBS-enabled setting, the Track B state includes a joint representation of classical and quantum-like components, coupled via a rank-1 tensor.

Let \(u_t \in \R^n\) be the classical feature vector (e.g., frequency-coded) and \(v_t \in \R^m\) be the quantum-like amplitude feature vector for Track B at time \(t\).
Let their suprema over the time horizon be
\[
  B_c = \sup_t \norm{u_t}_2, \qquad
  B_q = \sup_t \norm{v_t}_2.
\]

\begin{definition}[Coupling tensor]
The Track B coupling tensor at time \(t\) is
\[
  T_t^{(B)} = u_t v_t^\top \in \R^{n \times m}.
\]
\end{definition}

\begin{proposition}[Norms of a rank-1 coupling tensor]
For \(T_t^{(B)} = u_t v_t^\top\),
\[
  \norm{T_t^{(B)}}_F = \norm{u_t}_2 \norm{v_t}_2, \qquad
  \norm{T_t^{(B)}}_{2\to2} = \norm{u_t}_2 \norm{v_t}_2,
\]
so in particular
\[
  \norm{T_t^{(B)}}_F \le B_c B_q, \qquad
  \norm{T_t^{(B)}}_{2\to2} \le B_c B_q.
\]
\end{proposition}

\begin{proof}
For a rank-1 matrix \(T = u v^\top\), the singular values are \(\{\norm{u}_2 \norm{v}_2,0,\dots,0\}\).
Hence the Frobenius norm satisfies
\[
  \norm{T}_F^2 = \sum_{i,j} u_i^2 v_j^2 = \norm{u}_2^2 \norm{v}_2^2,
\]
so \(\norm{T}_F = \norm{u}_2 \norm{v}_2\).
The spectral norm (largest singular value) is \(\norm{u}_2 \norm{v}_2\).
The inequalities follow by the definitions of \(B_c\) and \(B_q\).
\end{proof}

\subsubsection{Track B Encryption Operator}

Let \(S_t^{(B)}\) denote an embedded Track B state vector in \(\R^d\) for some fixed \(d\), such that \(T_t^{(B)} S_t^{(B)}\) and \(\mathbf{M}_t^{(B)} T_t^{(B)} S_t^{(B)}\) are well-defined.
Let \(F_B : \R^d \to \R^d\) be a (possibly nonlinear) feedback function.

\begin{definition}[Track B encryption operator]
The Track B encryption operator at time \(t\) is
\[
  E_B(t) = \mathbf{M}_t^{(B)} \bigl( T_t^{(B)} S_t^{(B)} \bigr) + F_B\bigl(S_t^{(B)}\bigr).
\]
\end{definition}

Assume there exists \(D_S, D_F > 0\) such that
\[
  \norm{S_t^{(B)}}_2 \le D_S, \qquad
  \norm{F_B(S_t^{(B)})}_2 \le D_F
\]
for all \(t\) in the horizon of interest.

\begin{proposition}[Norm bound on \(E_B(t)\)]
Under the assumptions above,
\[
  \norm{E_B(t)}_2 \le N_t B_c B_q D_S + D_F.
\]
\end{proposition}

\begin{proof}
Using the triangle inequality and submultiplicativity of operator norms,
\begin{align*}
  \norm{E_B(t)}_2
  &= \norm{\mathbf{M}_t^{(B)} T_t^{(B)} S_t^{(B)} + F_B(S_t^{(B)})}_2 \\
  &\le \norm{\mathbf{M}_t^{(B)} T_t^{(B)} S_t^{(B)}}_2 + \norm{F_B(S_t^{(B)})}_2 \\
  &\le \norm{\mathbf{M}_t^{(B)}}_{2\to2} \norm{T_t^{(B)}}_{2\to2} \norm{S_t^{(B)}}_2 + D_F.
\end{align*}
By the earlier bounds,
\[
  \norm{\mathbf{M}_t^{(B)}}_{2\to2} \le N_t, \quad
  \norm{T_t^{(B)}}_{2\to2} \le B_c B_q, \quad
  \norm{S_t^{(B)}}_2 \le D_S,
\]
so
\[
  \norm{E_B(t)}_2 \le N_t B_c B_q D_S + D_F.
\]
\end{proof}

\subsection{HKDF, Transcript Binding, and Multiplicity Labels}

\subsubsection{Transcript Chains and Context Hash}

Let \(\SHA\) be a hash function modeled as a random oracle.
For each role \(s \in \{A,B\}\), define a per-role transcript chain:

\begin{definition}[Per-role transcript chain]
Initialize
\[
  \mathrm{chain}_{s,0} = 0^{256}.
\]
For \(t \ge 1\), given the \(t\)-th message \(m_{s,t}\) sent by role \(s\),
\[
  \mathrm{chain}_{s,t} = \SHA\bigl( \mathrm{chain}_{s,t-1} \,\|\, m_{s,t} \bigr).
\]
\end{definition}

\begin{definition}[Track B context hash]
The Track B context hash at time \(t\) is
\[
  \mathrm{ctx}_t^{(B)} = \SHA\bigl( \mathrm{chain}_{A,t} \,\|\, \mathrm{chain}_{B,t} \bigr).
\]
\end{definition}

These definitions mirror standard hash-chain transcript binding patterns used in secure messaging and zero-knowledge transcript binding.\cite{CommitmentSchemeWiki}

\subsubsection{Multiplicity Label Encoding}

\begin{definition}[Multiplicity label encoding]
Given multiplicity vector \(M_t^{(B)}\), define a fixed-length encoding
\[
  \mathrm{mult}_t^{(B)} = \mathrm{Enc}_8\bigl(M_t^{(B)}\bigr) \in \{0,1\}^{64},
\]
where \(\mathrm{Enc}_8\) is an injective (or collision-resistant) map from the relevant multiplicity range into 8 bytes.
\end{definition}

This label is treated as authenticated auxiliary information in the key-derivation process, not as a source of cryptographic entropy.

\subsubsection{HKDF-Based Directional Key Derivation}

Let \(\HKDF\) denote an HMAC-based key-derivation function following the extract-then-expand paradigm (e.g., RFC 5869).\cite{HKDFWiki,HKDFDotNet,GopherHKDF}

\begin{definition}[Directional info string]
For direction \(d \in \{A \to B, B \to A\}\),
\[
  \mathrm{info}_d^{(B)} =
    \texttt{"QSBS-v1:enc:"} \,\|\, d \,\|\, \mathrm{ctx}_t^{(B)} \,\|\, \mathrm{mult}_t^{(B)}.
\]
\end{definition}

\begin{definition}[Directional AEAD key for Track B]
Let \(K_t\) be a shared secret (e.g., from Diffie–Hellman or QKD) and \(\mathrm{salt}_t\) a salt.
The directional AEAD key at time \(t\) is
\[
  K_{\mathrm{enc},d,t}^{(B)} = \HKDF\bigl(K_t; \mathrm{salt}_t, \mathrm{info}_d^{(B)}\bigr).
\]
\end{definition}

\begin{remark}
Standard analyses of HKDF show that if \(K_t\) has sufficient min-entropy and HMAC behaves as a pseudorandom function, then HKDF outputs are indistinguishable from random keys.\cite{HKDFWiki,HKDFDotNet,GopherHKDF}
Including \(\mathrm{ctx}_t^{(B)}\) and \(\mathrm{mult}_t^{(B)}\) in \(\mathrm{info}_d^{(B)}\) binds the derived keys to a specific transcript and multiplicity label without degrading security, provided these values are public and not reused inconsistently.
\end{remark}

\subsection{AEAD Encryption and Associated Data}

Let \(\Enc_K(\mathrm{nonce}, m, \mathrm{ad})\) be an AEAD encryption algorithm, such as AES-GCM, with key \(K\), nonce, plaintext \(m\), and associated data \(\mathrm{ad}\).
Assume the AEAD is IND-CPA and INT-CTXT secure.

\begin{definition}[Track B associated data]
Let \(x_t^{(B)}\) be the Track B plaintext payload at time \(t\), and let \(F_t^{(B)}\) and \(M_t^{(B)}\) be the associated frequency vector and multiplicity operator.
Define associated data as the serialized tuple
\[
  \mathrm{ad}_t^{(B)} = \mathrm{enc}\bigl(\mathrm{ctx}_t^{(B)}, F_t^{(B)}, M_t^{(B)}\bigr).
\]
\end{definition}

\begin{definition}[Track B ciphertext]
Let \(\mathrm{nonce}_{d,t}\) be a nonce.
The Track B ciphertext for direction \(d\) at time \(t\) is
\[
  C_t^{(B)} = \Enc_{K_{\mathrm{enc},d,t}^{(B)}}\bigl(\mathrm{nonce}_{d,t}, x_t^{(B)}, \mathrm{ad}_t^{(B)}\bigr).
\]
\end{definition}

\begin{remark}
Because AEAD schemes authenticate the associated data, any attempt to tamper with \(\mathrm{ctx}_t^{(B)}\), \(F_t^{(B)}\), or \(M_t^{(B)}\) will lead to decryption failure.
This provides integrity and binding between the ciphertext, the transcript state, and the multiplicity operator.
\end{remark}

\subsubsection{Feedback Dynamics and Contraction}

\subsubsection{State Space and Metric}

Let the Track B feedback state be a pair
\[
  Z_t^{(B)} = \bigl(M_t^{(B)}, T_t^{(B)}\bigr) \in \R^k \times \R^{n \times m}.
\]
Equip this product space with the norm
\[
  \norm{Z_t^{(B)}} := \norm{M_t^{(B)}}_2 + \norm{T_t^{(B)}}_F.
\]

\begin{lemma}[Completeness]
The space \(\R^k \times \R^{n \times m}\) with the metric \(d(Z,Z') = \norm{Z-Z'}\) is a complete metric space.
\end{lemma}

\begin{proof}
Both \(\R^k\) and \(\R^{n \times m}\) are finite-dimensional normed spaces and hence complete; their Cartesian product equipped with the sum norm is also complete.
\end{proof}

\subsubsection{Feedback Map and Lipschitz Condition}

Let \(\Theta_t\) denote a vector of observed metrics at time \(t\) (e.g., error rates, latency, QBER, load).
Define a feedback map \(f_B\) via
\[
  Z_{t+1}^{(B)} = f_B\bigl(Z_t^{(B)}, \Theta_t\bigr).
\]

\begin{definition}[Lipschitz condition]
Suppose there exist non-negative constants \(L_M,L_T\) such that for all \(Z = (M,T)\) and \(Z' = (M',T')\),
\begin{align*}
  \norm{f_B(Z,\Theta)_M - f_B(Z',\Theta)_M}_2 &\le L_M \norm{Z-Z'}, \\
  \norm{f_B(Z,\Theta)_T - f_B(Z',\Theta)_T}_F &\le L_T \norm{Z-Z'}.
\end{align*}
Let \(L = L_M + L_T\).
\end{definition}

\begin{proposition}[Global Lipschitz bound]
Under the above conditions,
\[
  \norm{f_B(Z,\Theta) - f_B(Z',\Theta)} \le L \norm{Z - Z'}
\]
for all \(Z,Z'\).
\end{proposition}

\begin{proof}
By definition,
\[
  \norm{f_B(Z,\Theta) - f_B(Z',\Theta)}
  = \norm{f_B(Z,\Theta)_M - f_B(Z',\Theta)_M}_2
    + \norm{f_B(Z,\Theta)_T - f_B(Z',\Theta)_T}_F.
\]
Applying the component-wise Lipschitz bounds yields
\[
  \norm{f_B(Z,\Theta) - f_B(Z',\Theta)}
  \le (L_M + L_T) \norm{Z-Z'} = L \norm{Z-Z'}.
\]
\end{proof}

\begin{theorem}[Banach contraction for Track B feedback]
If \(L < 1\), then:
\begin{enumerate}
  \item There exists a unique fixed point \(Z_\ast^{(B)}\) such that
  \[
    f_B\bigl(Z_\ast^{(B)},\Theta\bigr) = Z_\ast^{(B)}
  \]
  for all admissible \(\Theta\) in the regime considered (e.g., steady-state metrics).
  \item For any initial state \(Z_0^{(B)}\), the iterations
  \[
    Z_{t+1}^{(B)} = f_B\bigl(Z_t^{(B)},\Theta_t\bigr)
  \]
  satisfy the contraction bound
  \[
    \norm{Z_t^{(B)} - Z_\ast^{(B)}} \le L^t \norm{Z_0^{(B)} - Z_\ast^{(B)}},
  \]
  provided the Lipschitz constants remain valid over the trajectory.
\end{enumerate}
\end{theorem}

\begin{proof}
By the previous lemma, the state space is complete.
If \(L<1\), then for fixed \(\Theta\) (or for \(\Theta_t\) staying in a regime where the same Lipschitz constants apply), the map \(Z \mapsto f_B(Z,\Theta)\) is a contraction.
The Banach Fixed-Point Theorem guarantees existence and uniqueness of \(Z_\ast^{(B)}\) and geometric convergence from any starting point.
\end{proof}

\subsubsection{Summary of Bounds and Guarantees}

For Track B, the key operator norm and security-relevant bounds can be summarized as follows:

\begin{itemize}
  \item \textbf{Multiplicity norms.}
  For \(M_t^{(B)} \in \N^k\),
  \[
    \norm{M_t^{(B)}}_1 = N_t, \quad
    \norm{M_t^{(B)}}_2 \le N_t, \quad
    \norm{M_t^{(B)}}_\infty \le N_t,
  \]
  and for the diagonal operator \(\mathbf{M}_t^{(B)}\),
  \(\norm{\mathbf{M}_t^{(B)}}_{2\to2} \le N_t\).
  \item \textbf{Coupling tensor.}
  For the rank-1 tensor \(T_t^{(B)} = u_t v_t^\top\),
  \[
    \norm{T_t^{(B)}}_F = \norm{u_t}_2 \norm{v_t}_2 \le B_c B_q, \quad
    \norm{T_t^{(B)}}_{2\to2} \le B_c B_q.
  \]
  \item \textbf{Encryption operator.}
  For bounded state \(\norm{S_t^{(B)}}_2 \le D_S\) and feedback \(\norm{F_B(S_t^{(B)})}_2 \le D_F\),
  \[
    \norm{E_B(t)}_2 \le N_t B_c B_q D_S + D_F.
  \]
  \item \textbf{HKDF security.}
  Provided \(K_t\) has sufficient min-entropy and HKDF is instantiated with a secure HMAC, \(K_{\mathrm{enc},d,t}^{(B)}\) is computationally indistinguishable from random, even when \(\mathrm{info}_d^{(B)}\) includes public multiplicity labels and context hashes.\cite{HKDFWiki,GopherHKDF}
  \item \textbf{AEAD binding.}
  AEAD security ensures integrity of both ciphertext and associated data; any change in \(\mathrm{ctx}_t^{(B)}\), \(F_t^{(B)}\), or \(M_t^{(B)}\) will cause decryption to fail, binding Track B ciphertexts to their specific transcript and multiplicity state.
  \item \textbf{Feedback contraction.}
  If the Track B feedback map satisfies the Lipschitz condition with \(L<1\), the multiplicity and coupling tensor converge to a unique fixed point at geometric rate, providing a mathematically controlled regime for adaptive key rotation or state evolution.
\end{itemize}

These results form the mathematical backbone of the Track B design and justify the use of prime-indexed multiplicities, rank-1 coupling tensors, and HKDF-based directional keys in the combined certification-and-transport architecture.

\vspace{1em}
\noindent\textbf{References (for background only)}
\begin{itemize}
  \item Standard references on Walsh--Hadamard transforms and orthogonal transforms.\cite{WHTOverview,WHTOrthogonal}
  \item Cryptographic references on Pedersen commitments, HKDF, and commitment schemes.\cite{PedersenZKDocs,PedersenRareSkills,CommitmentSchemeWiki,HKDFWiki,GopherHKDF}
\end{itemize}


\section{Track C: Signal Model and Windowed Feature Space}

\subsection{WESAD Multichannel Signals}

Track~C operates on multichannel time-series signals derived from the WESAD dataset (ECG, EDA, RESP, TEMP, ACC, etc.), but the mathematics below assumes a general multichannel sequence model.

\begin{definition}[Raw multichannel signal]
Let \(C \in \N\) be the number of sensor channels.
For a given subject \(s\), define the discrete-time multichannel signal
\[
  X_s(t) = \bigl(x_{s,c}(t)\bigr)_{c=1}^C \in \R^C, \quad t = 0,1,2,\dots .
\]
\end{definition}

\begin{definition}[Windowed segments]
Fix a window length \(L \in \N\) and stride \(R \in \N\).
The \(j\)-th window for subject \(s\) is the tensor
\[
  X_{s,j} = \bigl(x_{s,c}(t)\bigr)_{c=1,\dots,C;\, t = t_j,\dots,t_j+L-1}
  \in \R^{C \times L},
\]
where \(t_j = t_0 + jR\).
\end{definition}

Each window is associated with an affect label \(y_{s,j}\) (e.g.~baseline, stress, amusement), but the label is not needed for the norm bounds.

\subsubsection{Feature Extraction}

\begin{definition}[Per-window feature vector]
For each window \((s,j)\), a deterministic feature extractor
\[
  \mathcal{F} : \R^{C \times L} \to \R^d
\]
produces a feature vector
\[
  f_{s,j} = \mathcal{F}(X_{s,j}) \in \R^d.
\]
\end{definition}

\begin{remark}
In practice, \(\mathcal{F}\) concatenates statistics such as means, variances, frequency-band powers, and domain-specific features (e.g.~HRV from ECG, tonic/phasic components from EDA). The derivations below treat \(\mathcal{F}\) as an arbitrary, fixed map and place assumptions only on the resulting \(f_{s,j}\).
\end{remark}

\subsection{Channel Aggregation and Operator Norms}

\subsubsection{Channel Aggregation Matrix}

The CAS view groups the \(d\) features into \(K\) channels.

\begin{definition}[Aggregation matrix]
Let \(K \in \N\) be the number of CAS channels.
The channel aggregation matrix is
\[
  A \in \R^{K \times d},
\]
and the channel-intensity vector for window \((s,j)\) is
\[
  a_{s,j} = A f_{s,j} \in \R^K.
\]
\end{definition}

No special structure is required, but a common pattern is that each row of \(A\) selects and sums features corresponding to a particular modality or semantic grouping.

\begin{definition}[Operator norms]
For a matrix \(M \in \R^{m \times n}\),
\[
  \norm{M}_{2\to 2} = \sup_{x \ne 0} \frac{\norm{M x}_2}{\norm{x}_2}
\]
denotes the spectral (operator 2-) norm. We also use the Frobenius norm
\[
  \norm{M}_F = \biggl( \sum_{i,j} M_{ij}^2 \biggr)^{1/2}
\]
when convenient.
\end{definition}

\begin{proposition}[Channel norm bound]
Suppose \(\norm{A}_{2\to 2} \le B_A\) for some \(B_A > 0\).
Then for any feature vector \(f \in \R^d\),
\[
  \norm{A f}_2 \le B_A \norm{f}_2.
\]
In particular, for each window \((s,j)\),
\[
  \norm{a_{s,j}}_2 \le B_A \norm{f_{s,j}}_2.
\]
\end{proposition}

\begin{proof}
By the definition of the operator 2-norm,
\[
  \norm{A f}_2 \le \norm{A}_{2\to2} \norm{f}_2 \le B_A \norm{f}_2.
\]
Substituting \(f = f_{s,j}\) yields the second statement.
\end{proof}

\subsubsection{Specific Structures of \texorpdfstring{$A$}{A} and Norm Estimates}

If each feature participates in at most one channel and all nonzero entries of \(A\) are equal to 1, then we can bound \(\norm{A}_{2\to2}\) in terms of row sparsity.

\begin{proposition}[Row-sparsity bound]
Assume \(A \in \{0,1\}^{K \times d}\) and that each column has at most one nonzero entry.
Let
\[
  r_k = \#\{ \ell : A_{k,\ell} = 1 \}, \quad k=1,\dots,K,
\]
and let \(r_{\max} = \max_k r_k\).
Then
\[
  \norm{A}_{2\to2} \le \sqrt{r_{\max}}.
\]
\end{proposition}

\begin{proof}
Compute
\[
  (A A^\top)_{k,k'} = \sum_{\ell=1}^d A_{k,\ell} A_{k',\ell}.
\]
Because each column has at most one nonzero, \(A_{k,\ell} A_{k',\ell} = 0\) whenever \(k \ne k'\).
Thus \(A A^\top\) is diagonal with
\[
  (A A^\top)_{k,k} = \sum_{\ell=1}^d A_{k,\ell}^2 = r_k.
\]
The eigenvalues of \(A A^\top\) are therefore \(\{r_k\}_{k=1}^K\).
The spectral norm of \(A\) satisfies
\[
  \norm{A}_{2\to2}^2 = \lambda_{\max}(A A^\top) = r_{\max},
\]
so \(\norm{A}_{2\to2} \le \sqrt{r_{\max}}\).
\end{proof}

\begin{corollary}[Norm bound with row-sparsity]
If \(A\) satisfies the assumptions above, then for each window \((s,j)\),
\[
  \norm{a_{s,j}}_2 \le \sqrt{r_{\max}} \,\norm{f_{s,j}}_2.
\]
\end{corollary}

\subsection{Orthogonal WHT Mixing Layer}

\subsubsection{Walsh--Hadamard-Type Matrix}

Track~C uses an orthogonal mixing matrix \(H_K\), often chosen as a normalized Walsh--Hadamard matrix when \(K\) is a power of two.

\begin{definition}[Orthogonal mixing matrix]
Let \(H_K \in \R^{K \times K}\) be a real orthogonal matrix:
\[
  H_K H_K^\top = I_K = H_K^\top H_K.
\]
Given a channel vector \(a_{s,j}\), the mixed spectral vector is
\[
  h_{s,j} = H_K a_{s,j} \in \R^K.
\]
\end{definition}

\begin{proposition}[Norm preservation under \(H_K\)]
For any \(x \in \R^K\),
\[
  \norm{H_K x}_2 = \norm{x}_2.
\]
In particular, for each window \((s,j)\),
\[
  \norm{h_{s,j}}_2 = \norm{a_{s,j}}_2.
\]
\end{proposition}

\begin{proof}
Orthogonality implies
\[
  \norm{H_K x}_2^2 = (H_K x)^\top (H_K x)
  = x^\top H_K^\top H_K x
  = x^\top x = \norm{x}_2^2.
\]
Taking square roots yields \(\norm{H_K x}_2 = \norm{x}_2\).
Substituting \(x = a_{s,j}\) gives the claim about \(\norm{h_{s,j}}_2\).
\end{proof}

\subsubsection{Composite Linear Map and Operator Norms}

The overall linear map from feature space to spectral space is
\[
  M = H_K A : \R^d \to \R^K, \quad M f = H_K (A f).
\]

\begin{proposition}[Operator norm of \(M = H_K A\)]
The operator 2-norm of \(M\) equals that of \(A\):
\[
  \norm{M}_{2\to2} = \norm{A}_{2\to2}.
\]
\end{proposition}

\begin{proof}
Because \(H_K\) is orthogonal, \(\norm{H_K x}_2 = \norm{x}_2\) for all \(x\).
Thus for any \(f\),
\[
  \norm{M f}_2 = \norm{H_K A f}_2 = \norm{A f}_2.
\]
Taking suprema over nonzero \(f\) yields
\[
  \norm{M}_{2\to2} = \sup_{f \ne 0} \frac{\norm{M f}_2}{\norm{f}_2}
  = \sup_{f \ne 0} \frac{\norm{A f}_2}{\norm{f}_2}
  = \norm{A}_{2\to2}.
\]
\end{proof}

\begin{corollary}[Spectral norm bound for Track C]
If \(\norm{A}_{2\to2} \le B_A\), then for each window \((s,j)\),
\[
  \norm{h_{s,j}}_2 \le B_A \norm{f_{s,j}}_2.
\]
\end{corollary}

\begin{proof}
We have \(h_{s,j} = M f_{s,j}\).
Thus
\[
  \norm{h_{s,j}}_2 = \norm{M f_{s,j}}_2 \le \norm{M}_{2\to2} \norm{f_{s,j}}_2 = \norm{A}_{2\to2} \norm{f_{s,j}}_2 \le B_A \norm{f_{s,j}}_2.
\]
\end{proof}

\subsection{CAS Feature Vector and Boundedness}

\subsubsection{Definition of Track C CAS Feature}

For each window, Track~C defines a CAS feature vector as some function of \(a_{s,j}\) and \(h_{s,j}\).

\begin{definition}[CAS feature map]
Let \(\Psi : \R^K \times \R^K \to \R^D\) be a fixed (possibly nonlinear) map.
The Track~C CAS feature for window \((s,j)\) is
\[
  z_{s,j} = \Phi(f_{s,j}) := \Psi\bigl(a_{s,j}, h_{s,j}\bigr)
  = \Psi\bigl(A f_{s,j}, H_K A f_{s,j}\bigr) \in \R^D.
\]
\end{definition}

Examples include simple concatenation \(\Psi(a,h) = (a,h)\), residuals, or normalized combinations.

\begin{definition}[Lipschitz constant of \(\Psi\)]
Suppose there exists \(L_\Psi > 0\) such that for all \((a,h)\) and \((a',h')\),
\[
  \norm{\Psi(a,h) - \Psi(a',h')}_2
  \le L_\Psi \bigl( \norm{a-a'}_2 + \norm{h-h'}_2 \bigr).
\]
\end{definition}

\begin{proposition}[Lipschitz constant of \(\Phi\)]
Assume \(\norm{A}_{2\to2} \le B_A\) and that \(\Psi\) is Lipschitz with constant \(L_\Psi\) as above.
Then the composite map \(\Phi : \R^d \to \R^D\) satisfies
\[
  \norm{\Phi(f) - \Phi(f')}_2 \le 2 B_A L_\Psi \norm{f - f'}_2
\]
for all \(f,f' \in \R^d\).
\end{proposition}

\begin{proof}
Write \(a = A f\), \(h = H_K A f\), and similarly for \(a',h'\).
Since \(H_K\) is orthogonal,
\[
  \norm{h - h'}_2 = \norm{H_K(A f - A f')}_2 = \norm{A(f-f')}_2 \le B_A \norm{f-f'}_2,
\]
and also \(\norm{a - a'}_2 = \norm{A(f-f')}_2 \le B_A \norm{f-f'}_2\).
Therefore,
\begin{align*}
  \norm{\Phi(f) - \Phi(f')}_2
  &= \norm{\Psi(a,h) - \Psi(a',h')}_2 \\
  &\le L_\Psi \bigl( \norm{a-a'}_2 + \norm{h-h'}_2 \bigr) \\
  &\le L_\Psi (B_A \norm{f-f'}_2 + B_A \norm{f-f'}_2) \\
  &= 2 B_A L_\Psi \norm{f-f'}_2.
\end{align*}
\end{proof}

\begin{corollary}[Uniform bound given bounded feature norm]
If \(\norm{f_{s,j}}_2 \le B_f\) for all \((s,j)\), then
\[
  \norm{a_{s,j}}_2 \le B_A B_f, \quad
  \norm{h_{s,j}}_2 \le B_A B_f,
\]
and, in particular, if \(\Psi\) is Lipschitz and \(\Psi(0,0) = 0\),
\[
  \norm{z_{s,j}}_2 \le 2 B_A L_\Psi B_f.
\]
\end{corollary}

\subsection{Logistic Regression on CAS Features}

\subsubsection{Model Definition}

For binary stress vs.\ non-stress, each window \((s,j)\) is labeled \(y_{s,j} \in \{0,1\}\).

\begin{definition}[Logistic regression on CAS features]
Let \(\theta \in \R^D\) be model parameters.
Define
\[
  p_\theta(y_{s,j}=1 \mid z_{s,j}) = \sigma(\theta^\top z_{s,j}),
\]
where \(\sigma(z) = \frac{1}{1+e^{-z}}\).
The per-sample negative log-likelihood is
\[
  \ell(\theta; z_{s,j}, y_{s,j})
  = -y_{s,j} \log \sigma(\theta^\top z_{s,j})
    - (1-y_{s,j}) \log \bigl(1-\sigma(\theta^\top z_{s,j})\bigr).
\]
\end{definition}

\subsubsection{Gradient Norm Bound}

\begin{proposition}[Gradient bound]
If \(\norm{z_{s,j}}_2 \le B_Z\) for all \((s,j)\), then for all \(\theta\),
\[
  \norm{\nabla_\theta \ell(\theta; z_{s,j}, y_{s,j})}_2 \le B_Z.
\]
\end{proposition}

\begin{proof}
The gradient is
\[
  \nabla_\theta \ell(\theta; z_{s,j}, y_{s,j})
  = \bigl(\sigma(\theta^\top z_{s,j}) - y_{s,j}\bigr) z_{s,j}.
\]
Since \(y_{s,j} \in \{0,1\}\) and \(\sigma(\theta^\top z_{s,j}) \in (0,1)\), we have
\(\abs{\sigma(\theta^\top z_{s,j}) - y_{s,j}} \le 1\).
Therefore,
\[
  \norm{\nabla_\theta \ell}_2
  \le \abs{\sigma(\theta^\top z_{s,j}) - y_{s,j}} \,\norm{z_{s,j}}_2
  \le \norm{z_{s,j}}_2 \le B_Z.
\]
\end{proof}

\begin{corollary}[Lipschitz gradient]
If \(B_Z < \infty\), the logistic loss has gradients uniformly bounded in norm by \(B_Z\), which facilitates step-size selection and convergence guarantees for first-order optimization methods.
\end{corollary}

subsection{Within-Subject Normalization and Baseline Deviations}

Track~C can also use within-subject baseline normalization for CAS features.

\subsubsection{Subject Baseline CAS Feature}

\begin{definition}[Subject baseline CAS feature]
For subject \(s\), let \(\mathcal{B}_s\) denote the set of baseline windows.
Define the baseline CAS mean as
\[
  \bar{z}_s^{(\mathrm{base})}
  = \frac{1}{\abs{\mathcal{B}_s}} \sum_{j \in \mathcal{B}_s} z_{s,j}.
\]
\end{definition}

\begin{definition}[Deviation from baseline]
For any window \((s,j)\), define the deviation
\[
  \Delta z_{s,j} = z_{s,j} - \bar{z}_s^{(\mathrm{base})}.
\]
\end{definition}

\subsubsection{Norm Bounds for Deviations}

\begin{proposition}[Deviation norm bound]
If \(\norm{z_{s,j}}_2 \le B_Z\) for all windows of subject \(s\), then
\[
  \norm{\Delta z_{s,j}}_2 \le 2 B_Z
\]
for all \(j\).
\end{proposition}

\begin{proof}
Using the triangle inequality,
\[
  \norm{\Delta z_{s,j}}_2
  = \norm{z_{s,j} - \bar{z}_s^{(\mathrm{base})}}_2
  \le \norm{z_{s,j}}_2 + \norm{\bar{z}_s^{(\mathrm{base})}}_2.
\]
Moreover,
\[
  \norm{\bar{z}_s^{(\mathrm{base})}}_2
  = \biggl\lVert \frac{1}{\abs{\mathcal{B}_s}} \sum_{j \in \mathcal{B}_s} z_{s,j} \biggr\rVert_2
  \le \frac{1}{\abs{\mathcal{B}_s}} \sum_{j \in \mathcal{B}_s} \norm{z_{s,j}}_2
  \le B_Z.
\]
Hence \(\norm{\Delta z_{s,j}}_2 \le B_Z + B_Z = 2 B_Z\).
\end{proof}

\subsection{Summary of Operator Norm Bounds}

For the Track~C WESAD CAS proof-of-concept, the key bounds are:

\begin{itemize}
  \item \textbf{Aggregation matrix:}
    \[
      \norm{A}_{2\to2} \le B_A, \quad
      \norm{a_{s,j}}_2 \le B_A \norm{f_{s,j}}_2.
    \]
    If \(A\) is a 0--1 matrix with each column used at most once and row sparsity \(r_{\max}\), then \(\norm{A}_{2\to2} \le \sqrt{r_{\max}}\).
  \item \textbf{Orthogonal mixing:}
    \[
      \norm{H_K x}_2 = \norm{x}_2, \quad
      \norm{h_{s,j}}_2 = \norm{a_{s,j}}_2.
    \]
  \item \textbf{Composite map:}
    \[
      M = H_K A, \quad
      \norm{M}_{2\to2} = \norm{A}_{2\to2}.
    \]
  \item \textbf{CAS feature map:}
    If \(\Psi\) is Lipschitz with constant \(L_\Psi\) and \(\norm{A}_{2\to2} \le B_A\), then
    \[
      \norm{\Phi(f) - \Phi(f')}_2 \le 2 B_A L_\Psi \norm{f-f'}_2.
    \]
    With bounded feature norm \(\norm{f_{s,j}}_2 \le B_f\), all CAS features satisfy \(\norm{z_{s,j}}_2 \le 2 B_A L_\Psi B_f\).
  \item \textbf{Logistic regression gradient:}
    If \(\norm{z_{s,j}}_2 \le B_Z\), then
    \[
      \norm{\nabla_\theta \ell(\theta; z_{s,j}, y_{s,j})}_2 \le B_Z
    \]
    for all \(\theta\) and all windows.
  \item \textbf{Baseline deviations:}
    With \(\norm{z_{s,j}}_2 \le B_Z\), the within-subject deviations satisfy
    \[
      \norm{\Delta z_{s,j}}_2 \le 2 B_Z.
    \]
\end{itemize}

These results collectively provide the explicit proofs and norm bounds underpinning Track~C’s WESAD CAS formulation, ensuring that the spectral and classifier layers are numerically stable and analytically well-controlled.



\nocite{*}
\bibliographystyle{unsrt}  
\bibliography{references}  


\end{document}
